% concatenated from main.tex
\documentclass[12pt,reqno]{amsart}

\usepackage{amsmath, amssymb, amsfonts, amsbsy, amsthm, latexsym, epsfig, bbm, amscd,graphicx}
\usepackage{dirtytalk}
\usepackage{comment}
\usepackage{xcolor}
\textwidth      6.5in
\oddsidemargin  0.0in
\evensidemargin 0.0in
\setlength{\topmargin}{.0in}
\setlength{\textheight}{8.5in}

\theoremstyle{definition}
\newtheorem{theorem}{Theorem} [section]
\newtheorem{conjecture}[theorem]{Conjecture}
\newtheorem{example}[theorem]{Example}
\theoremstyle{definition}
\newtheorem{corollary}[theorem]{Corollary}
\newtheorem{lemma}[theorem]{Lemma}
\newtheorem{proposition}[theorem]{Proposition}
\newtheorem{definition}[theorem]{Definition}
\newtheorem{exercise}[theorem]{Exercise}
\newtheorem{notation}[theorem]{Notation}
\newtheorem{remark}[theorem]{Remark}
\newtheorem*{namedtheorem}{Theorem}
\newtheorem*{namedlemma}{Lemma}
\newtheorem*{namedproposition}{Proposition}
\newtheorem*{namedcorollary}{Corollary}
\numberwithin{equation}{section}
\definecolor{vbcolor}{rgb}{0.35, 0.15, 0.6}
\newcommand{\vb}[1]{\textcolor{vbcolor}{#1}}

\begin{document}
%working title.
\title{Submodules of $H^2(\mathbb{T}^2)$ and frames by pairs of bounded commuting operators}

\author{Victor Bailey and Carlos Cabrelli}

\address{Department of Mathematics,
University of Oklahoma,
Norman, OK 73019 USA}
\email{victor.bailey@ou.edu}


\address{Departmentamento de Matemática, Universidad de Buenos Aires, Instituto de Matemática ``Luis Santaló'' (IMAS-CONICET-UBA), Buenos Aires, Argentina}
\email{cabrelli@dm.uba.ar}

\begin{abstract}
Recent work in Dynamical Sampling has been centered on characterizing frames  obtained by the orbit of a vector under a bounded operator. We prove a necessary and sufficient condition for a pair of bounded commuting operators on a separable infinite-dimensional Hilbert space to generate a frame by unilateral iterations on a single vector. Applying the theory on submodules of the Hardy module $H^2(\mathbb{T}^2)$, we characterize these frames in terms of their relation to the two-variable Jordan block on a certain quotient module and provide some properties of frames of this form. 
\end{abstract}
\maketitle
\section
{Introduction}


\quad A central problem in Dynamical Sampling  \cite{ACMT17} asks for the conditions needed to obtain frames and/or bases for a separable Hilbert space $H$ from collections of vectors of the form $\{A^n g\}_{{0 \leq n \leq \Psi (g)} , \, g \, \in \, G}$ 
where $A \in B(H)$, $G$ is a countable subset of $H,$ and $\Psi$ is a map from $G$ to $\mathbb{N}_0 \cup \{\infty\}$ (where $\mathbb{N}_0 = \mathbb{N} \cup \{0\}$).  As a complete solution to this Dynamical Sampling problem in finite-dimensional Hilbert spaces was provided in \cite{ACMT17}, recent works consider the problem where $H$ is a separable infinite-dimensional Hilbert space. In this setting, typically the problem is posed in the following way: Given an operator $T \in B(H)$ and a vector $\varphi \in H$, what are conditions for the system $\{T^n \varphi\}_{n \geq 0}$ to be a frame or have properties that are useful for reconstructing $H$ such as completeness, exactness, etc.

 \quad A necessary and sufficient condition for a system $\{T^n \varphi\}_{n \geq 0} \subset H$ to be a frame was given in \cite{CHP20}. This result exhibited the inextricable connection between Dynamical Sampling and Hardy Space theory as the characterization of the bounded operators that can be used to generate a frame for $H$ hinges on the relationship between the iterating operator and a compression of the shift operator to a certain type of subspace (known as a model space) of the Hardy Space, \begin{displaymath}
H^2(\mathbb{T}) = \{f \in L^2(\mathbb{T}) : \int_{\mathbb{T}} f(z) \overline{z}^n dz = 0 \, \,  \forall \, n<0 \}.
 \end{displaymath}
 %By Beurling's theorem, every invariant subspace of the shift operator, $S$, on $H^2(\mathbb{T})$ has the form $\theta H^2(\mathbb{T})$ where $\theta(\xi) \in H^2(\mathbb{T})$ is an inner function. Shift-invariant subspaces of $\ell^2(\mathbb{N}_0)$ correspond to shift-invariant subspaces of $H^2$ in a natural way due to the unitary equivalence of the operators $S$ and $R: \ell^2(\mathbb{N}_0) \to \ell^2(\mathbb{N}_0)$ where $S(x_0, x_1, \ldots) = (0, x_0, x_1, \ldots)$. 
 In Theorem 3.6 of \cite{CHP20} it is shown that for any complex separable infinite-dimensional Hilbert Space $H$, the collection of vectors $\{T^{n}\varphi\}_{n \geq 0} \subset H$ is a frame for $H$ if and only if for some inner function $\theta \in H^2(\mathbb{T})$ there is an infinite-dimensional model space $K_\theta =  H^2(\mathbb{T}) \ominus \theta H^2(\mathbb{T})$ for which the operator $T$ is similar to the compressed shift operator $S_{\theta} = P_{K_{\theta}} \left.S\right|_{K_{\theta}}$ (where $P_{K_{\theta}}$ the orthogonal projection onto $K_{\theta}$) and $\varphi$ is the image of the function $P_{K_{\theta}}1_{\mathbb{T}}$ under the invertible map given in the similarity relation between $T$ and $S_{\theta}$. %In fact, the inner function $\theta$ is unique and the model space $K_{\theta} = H^2(\mathbb{T}) \ominus Ker(V)$ with $V = U \mathcal{F}$ where $U$ is the synthesis operator of the frame $\{T^{n}\varphi\}_{n \in \mathbb{N}_0}$ and $\mathcal{F}$ is the Fourier transform. Letting $W = \left.V\right|_{K_{\theta}}$, then we have $T = WS_{\theta}W^{-1}$ and $\varphi =WP_{K_{\theta}}1_{\mathbb{T}}$. 

%equivalence of $H^2(\mathbb{T}^2)$ and $$H^2(\mathbb{D}^2)$ given in Yang operator theory bidisc I.

Recall that a system $\{f_n\}_{n \in I} \subset H$ is said to be a frame if there exist fixed constants $ 0 < C_{1} \leq C_{2}$ such that for each $ f \in H$,  
\begin{displaymath} C_{1} \|f\|^2 \leq \sum_{n \in I} | \langle f , f_n \rangle |^2 \leq C_{2} \|f\|^2 . \quad 
\end{displaymath}
A frame for $H$ admits basis-like expansions for each element of $H$. That is, for a system $\{f_n\}_{n \in I} \subset H $, that satisfies the inequality above, each vector $f\in H$ can be given a series representation of the form  $f = \underset{n \in I}{\sum} c_n \, f_n $
where $\{c_n\}_{n \in I} \in \ell^2(I)$ and convergence is in the norm of $H$.  For a frame, these expansions are not necessarily unique so that a frame can be viewed as a generalized basis for $H$. Frames for which these expansions are not unique, so that the frame is not minimal, are known as overcomplete frames and a minimal frame is known as a Riesz basis. 
%The utility of such a collection of vectors is made apparent from the fact that these vectors allow for the reconstruction of every element of $H$ as a series in terms of the vectors in the set, like a basis, while not necessarily satisfying the stringent requirements needed to be a basis. In fact, frames are simply collections of vectors that satisfy a particular inequality, known as the frame condition. 

%Frames were first introduced in the foundational paper \cite{SD52} by Duffin and Schaeffer and have applications in various areas of mathematics [cite some books on frames with applications].  

 In the sequel, we show that for frames of the form $\{T_{1}^iT_{2}^j\varphi\}_{i,j \geq 0}$ (where $T_{1}, T_{2} \in B(H)$ commute) an analogous relation holds as the one given in Theorem 3.6 of \cite{CHP20}. That is, the bounded and commuting iterating operators $T_{1}$ and $T_{2}$ must be similar to compressions of the two shift operators defined on a certain subspace of the Hardy Space on the bidimensional torus, $H^2(\mathbb{T}^2)$. 
In \cite{ACCP22}, the authors investigated a mixed setting in which one of the operators is applied bilaterally, while the other is applied only unilaterally. This asymmetric structure requires the development of specific techniques tailored to that framework. These methods are substantially different from those employed in our symmetric setting, which introduces its own significant analytical challenges, most notably, the need to work within the Hardy space on the bidisc.
In the space 
\begin{displaymath}
H^2(\mathbb{T}^2) = \{f(z, w) \in L^2(\mathbb{T}^2): \int_{\mathbb{T}^2} f(z,w) \overline{z}^m \overline{w}^n d \mu = 0 \, \, \text{if} \, \, \,  m<0 \, \, \text{or} \, \, n<0 \} 
\end{displaymath}
where $\mu$ is the Haar measure on $\mathbb{T}^2$, many of the characteristics of the space $H^2(\mathbb{T})$ are reflected. On $H^2(\mathbb{T}^2)$ we have two shift operators defined by multiplication by $z$ and multiplication by $w$ respectively and say a subspace $M \subseteq H^2(\mathbb{T}^2)$  is shift-invariant if it is invariant under both shift operators. However, the subspaces corresponding to the model space $K_\theta$ in this setting are not necessarily of the form $H^2(\mathbb{T}^2) \ominus \phi H^2(\mathbb{T}^2)$ as Beurling's characterization of shift-invariant subspaces of $H^2(\mathbb{T})$ does not carry over fully to this setting. For this reason, in order to seamlessly extend the result from Theorem 3.6 of \cite{CHP20} to the operator pair case, we make use of the theory on the structure of shift-invariant subspaces of multidimensional Hardy Spaces which enables us to obtain Corollary \ref{B-type} which follows from the  complete characterization of frames of the form $\{T_{1}^iT_{2}^j\varphi\}_{i,j \geq 0}$ where $T_{1}, T_{2} \in B(H)$ commute that we provide in Theorem \ref{char} below.
\begin{comment}Beurling's characterization of invariant subspaces does not translate fully to $H^2(\mathbb{T}^2)$, many works [...] characterize these subspaces in special cases.  
Note that in Theorem 6 [Christensen Philipp] it was necessary that the inner function $\varphi(\xi)$ for which $Ker(V) = \varphi H^2(\mathbb{T})$ was not a finite Blaschke product. By Theorem 3.14 [Invar Sub Radjavi], an invariant subspace $\varphi H^2(\mathbb{T})$ has finite codimension if and only if the inner function, $\varphi$, is a finite Blaschke product. Thus in Theorem 6, it was necessary that $\varphi H^2(\mathbb{T})$ (or in their notation $hL^{2}_+)$ had infinite codimension.  By [Argawal, Clark, Douglas] when  $\phi(z, w)$ is an inner function, any invariant subspace $M \subseteq H^2(\mathbb{T}^2)$ satisfies $ \phi M \subseteq M$ with equality if and only if  $\phi$ is constant. As stated in section 2 page 4 as well as in the proof of Corollary 3 [Argawal, Clark, Douglas], an invariant subspace $M \subseteq H^2(\mathbb{T}^2)$ with finite codimension has full range (definition on pg 5 of argawal douglas clark). As mentioned in [mandrekar], invariant subspaces of the form $\phi H^2(\mathbb{T}^2)$ do not have full range unless $\phi$ is constant. Thus, a Beurling type invariant subspace, $\phi H^2(\mathbb{T}^2)$, cannot have finite codimension unless $\phi$ constant. That is, invariant subspaces of the form $\phi H^2(\mathbb{T}^2)$ with $\phi(z, w)$ inner,  have infinite codimension unless $\phi H^2(\mathbb{T}^2) = H^2(\mathbb{T}^2)$. 
\end{comment}
\section{Shift-Invariant Subspaces of $H^2(\mathbb{T}^2)$}
%Can we show that every invariant subspace of $H^2(\mathbb{T}^2)$ that does not have full range is of Beurling-Type?
Note that the Hardy space has an equivalent definition as a collection of functions defined on the open unit disc $\mathbb{D}$. That is, the space \begin{displaymath}
    H^2(\mathbb{D}) = \{ f: \mathbb{D} \to \mathbb{C} \, \, \text{such that} \, \, f(z) = \underset{n \geq 0}{\sum} c_{n}z^n \, \, \text{where} \, \, \, (c_{n}) \in \ell^2(\mathbb{N}_0)\},
\end{displaymath} 
of analytic functions on $\mathbb{D}$ is known as the Hardy space on the disc and we identify this space of functions with those in $H^2(\mathbb{T})$ via radial limits \cite{MR07}. 
We say a subspace $M \subseteq H^2(\mathbb{T})$  is shift-invariant if it is invariant under the multiplication operator $S \in B(H^2(\mathbb{T}))$ where $Sf(z) = zf(z)$ for all $f \in H^2(\mathbb{T})$.

Similarly, we can identify functions in 
\begin{displaymath}
    H^2(\mathbb{D}^2) = \{ f: \mathbb{D}^2 \to \mathbb{C} \, \, \text{such that} \, \, f(z, w) = \underset{n \geq 0}{\sum} c_{n,m}z^nw^m \, \, \text{where} \, \, \, (c_{n, m}) \in \ell^2(\mathbb{N}_0\times \mathbb{N}_0)\},
\end{displaymath} 
the Hardy Space on the bidisc with those in $H^2(\mathbb{T}^2)$ via radial limits in like fashion. 

 We say a subspace $M \subseteq H^2(\mathbb{T}^2)$  is shift-invariant if it is invariant under both shifts and such subspaces are also referred to in the literature as submodules of the Hardy module $H^2(\mathbb{T}^2)$ \cite{CG03,S15,Y18}.

\begin{definition}
A subspace $M \subseteq H^2(\mathbb{T}^2)$  is shift-invariant, or a submodule of $H^2(\mathbb{T}^2)$, if it is invariant under both shift operators $S_z$ and $S_w$. That is, $M$ is shift-invariant if $S_zM = zM \subset M$ and $S_wM = wM \subset M$.
\end{definition}
 \begin{definition}
     A quotient module is a subspace in $H^2(\mathbb{T}^2)$ of the form $K = H^2(\mathbb{T}^2) \ominus M$ for some submodule $M \subset H^2(\mathbb{T}^2)$.
 \end{definition}
For a given submodule $M$ we associate the pair $(S_{K_z}, S_{K_w})$ to the quotient module $K= H^2(\mathbb{T}^2) \ominus M$ which are each compressions of the shift operators $S_z$ and $S_w$ to the quotient module. In the literature, this pair is known as a two-variable Jordan block \cite{LYY11}.
 \begin{definition}
For a given quotient module $K = H^2(\mathbb{T}^2) \ominus M$ the pair $(S_{K_z}, S_{K_w})$ where $S_{K_z} = P_{K} \left.S_{z}\right|_{K} $ and $S_{K_w} = P_{K} \left.S_{w}\right|_{K}$ (where $P_{K}$ is the orthogonal projection onto $K$) is the two-variable Jordan block on the quotient module $K$.
\end{definition}
Over the past several decades \cite{ACD86,CG03,R69,Y18} the structure of shift-invariant subspaces in $H^2(\mathbb{T}^2)$ ( or in $H^2(\mathbb{D}^2)$) has been the subject of extensive study in the area of multivariable operator theory. This has led to the formulation of the viewpoint of such subspaces as Hilbert modules \cite{CG03,DP89,S15} as subspaces closed under multiplication by $z \, \, \text{and} \, \, w$ are in fact closed under multiplication by all functions in the bidisc algebra $A(\mathbb{D}^2)$ (which is the closure of the polynomials in $C(\overline{\mathbb{D}}^2)$), and to the construction of various tools which can be used to ascertain the properties (such as the rank, codimension, etc.) of a given submodule \cite{Y18}. 
\newline
In the single variable case, that is in the space $H^2(\mathbb{T})$, there is a complete characterization of the structure of shift-invariant subspaces due to Beurling in terms of the inner functions which parametrize the spaces.  
  \begin{definition}
      An inner function in $H^2(\mathbb{D})$ is a function $\theta$ that satisfies $| \theta (z)| = 1$ almost everywhere on $\mathbb{T}$.
\end{definition}
Inner functions in $H^2(\mathbb{D})$ (or $H^2(\mathbb{T})$) are the class of unimodular functions in $H^2(\mathbb{D})$.
\begin{definition}
    A shift-invariant subspace $M$ is said to be unitarily equivalent to $N$ if there exists a unimodular function $\psi \in L^\infty$ such that $M =\psi N $.
\end{definition}
 \quad \, The following theorem characterizing the structure of the shift-invariant subspaces in $H^2(\mathbb{T})$ by Beurling \cite{B48} shows that they  are all unitarily equivalent to the full space $H^2(\mathbb{T})$.
 
   \begin{theorem}[Beurling]
    Every nontrivial invariant subspace of the shift operator $S \in B(H^2(\mathbb{T})) $, where $Sf(z) = zf(z)$, is of the form $\theta H^2(\mathbb{T})$ for some inner function $\theta$. Conversely, any  subspace of the form $\theta H^2(\mathbb{T})$ for some inner function $\theta$ is shift-invariant.
    \end{theorem}

There are several known examples of shift-invariant subspaces in $H^2(\mathbb{T}^2)$ that are not of the form given in Beurling's Theorem. That is, there exist submodules in $H^2(\mathbb{T}^2)$ that are not unitarily equivalent to $H^2(\mathbb{T}^2)$. For this reason numerous works have been aimed at providing characterizations for the structure of submodules \cite{Y18}.  




\begin{definition}
  An inner function in $H^2(\mathbb{T}^2)$ is a function $\phi$ that satisfies $| \phi (z, w)| = 1$ almost everywhere on $\mathbb{T}^2$.
\end{definition}

As noted above, Beurling's characterization of shift-invariant subspaces does not translate fully to the multidimensional Hardy Space setting as there exist shift-invariant subspaces in $H^2(\mathbb{T}^2)$ that are not of Beurling-type \cite{R69}, that is of the form $\phi H^2(\mathbb{T}^2)$ for some inner function $\phi \in H^2(\mathbb{T}^2)$.  However, due to Mandrekar \cite{M88}, we have a characterization of the Beurling-type shift-invariant subspaces in $H^2(\mathbb{T}^2)$.

\begin{definition}
A pair of isometries $V_1$ and $V_2$ doubly commute on a Hilbert space, $H$, if $V_1$ commutes with $V_2$ and $V_1$ commutes with $V_2^{*}$.
\end{definition}

%The following theorem characterizes the Beurling-type shift-invariant subspaces of $H^2(\mathbb{T}^2)$.
\begin{theorem}[Mandrekar]
A nontrivial shift-invariant subspace $M \subset H^2(\mathbb{T}^2)$ is of the form $\phi H^2(\mathbb{T}^2)$ with $\phi(z,w)$ inner if and only if $S_z$ and $S_w$ doubly commute on $M$.
 \end{theorem}

It is important to note that a nontrivial Beurling-type shift-invariant subspace, $\phi H^2(\mathbb{T}^2)$,  always has the property that the associated quotient module is infinite dimensional (that is, $dim(H^2(\mathbb{T}^2) \ominus \phi H^2(\mathbb{T}^2)) = \infty$). In Theorem 3.6 in \cite{CHP20} it was required that the subspace $K_\theta = H^2(\mathbb{T}) \ominus \theta H^2(\mathbb{T})$ had infinite dimension (which equivalently says $\theta H^2(\mathbb{T})$ has infinite codimension), since the theorem necessitates that the inner function $\theta$ not be a finite Blaschke product (which by Theorem 3.14 \cite{RR03} means $\theta H^2(\mathbb{T})$ has infinite codimension). 
Thus, the value of a Beurling-type submodule in $H^2(\mathbb{T}^2)$, for our purposes, is exhibited by the fact that we do not need to impose any conditions on the inner function in this setting to obtain $\phi H^2(\mathbb{T}^2)$ has infinite codimension. That is, we automatically satisfy the condition that the quotient module corresponding to $K_\theta$ in our setting has infinite dimension with any proper Beurling-type shift-invariant subspace. %Thus in Theorem 3.6, it was necessary that $\varphi H^2(\mathbb{T})$ had infinite codimension. 

\begin{proposition}\label{infcodim}
Every proper submodule of the form $\phi H^2(\mathbb{T}^2)$, where $\phi(z, w)$ is a (nonconstant) inner function, satisfies the property that $dim(H^2(\mathbb{T}^2) \ominus \phi H^2(\mathbb{T}^2)) = \infty$. That is, every proper Beurling-type shift-invariant subspace of $H^2(\mathbb{T}^2)$ has infinite codimension.
\end{proposition}
\begin{proof}
By \cite{ACD86}, a shift-invariant subspace $M \subseteq H^2(\mathbb{T}^2)$ with finite codimension has full range (definition on pg 5 of \cite{ACD86}). By \cite{M88}, shift-invariant subspaces of the form $\phi H^2(\mathbb{T}^2)$, where $\phi(z, w)\in H^2(\mathbb{T}^2)$ is an inner function, do not have full range unless $\phi$ is constant. Therefore, a Beurling-type shift-invariant subspace, $\phi H^2(\mathbb{T}^2)$, cannot have finite codimension unless $\phi$ is constant. That is, shift-invariant subspaces of the form $\phi H^2(\mathbb{T}^2)$  have infinite codimension unless $\phi H^2(\mathbb{T}^2) = H^2(\mathbb{T}^2)$ (since $\phi H^2(\mathbb{T}^2) = H^2(\mathbb{T}^2)$ if and only if $\phi$ is constant).
\end{proof}
A characterization of the submodules $M$ of finite codimension was provided by Ahern and Clark in terms of the topological properties of the zero set of the polynomials in $M$ \cite{Y18}.
\begin{comment}
All invariant subspaces with finite codimension have full range. There exist full range invariant subspaces that have infinite codimension. Each Beurling-type invariant subspace does not have full range. Are the invariant subspaces that do not have full range precisely those of Beurling-type \cite{ACD86}?
\end{comment}

\begin{comment}
For each of the following subspaces, the closure is taken w.r.t. the $L^{2}$-$norm$.
\begin{definition}
$i.)$ $\mathcal{H}_1 = \overline{span} \{z^m, w^n : m \geq 0, n \in \mathbb{Z}\}$. 
\newline
$ii.)$ $\mathcal{H}_2 = \overline{span} \{z^m, w^n : m \in \mathbb{Z}, n \geq 0 \}$.
\end{definition}

\begin{definition}
Let $M \subset H^2(\mathbb{T}^2)$ be shift-invariant, $L_z^{\infty}$ the weak$^*$-closure of the algebra generated by $\{1_{\mathbb{T}^2}, z, \overline{z}\}$, and $L_w^{\infty}$ the weak$^*$-closure of the algebra generated by $\{1_{\mathbb{T}^2}, w, \overline{w}\}$.
\newline
$i.)$ $(M)_2 = \overline{span} \{L_z^{\infty} M\}$
\newline
$ii.)$ $(M)_1 = \overline{span} \{L_w^{\infty} M\}$
\end{definition}

By Theorem 13 in \cite{GM88} we have that every shift-invariant subspace of $ M \subset H^2(\mathbb{T}^2)$ that is not full range satisfies either $(M)_1 \subsetneq \mathcal{H}_1$ or  $(M)_2 \subsetneq \mathcal{H}_2$.
\newline
From Lemma 12 in \cite{GM88} we have that  $(M)_2 = \overline{span}
\{ \underset{n \geq 0}{\bigcup}{\overline{z}^n M} \}$ and similarly $(M)_1 = \overline{span}
\{ \underset{n \geq 0}{\bigcup}{\overline{w}^n M} \}$.

The question posed in \cite{ACD86} is whether every invariant subspace, $M$, that does not have full range can be written in the form $\phi N$ for some invariant subspace $N \subset H^2(\mathbb{T}^2)$. We should try to show that if $\overline{span}
\{ \underset{n \geq 0}{\bigcup}{\overline{z}^n M} \} \subsetneq \mathcal{H}_2$ or  $\overline{span}
\{ \underset{n \geq 0}{\bigcup}{\overline{w}^n M} \} \subsetneq \mathcal{H}_1$, then there exists an inner function $\phi$ and some invariant subspace $N \subset H^2(\mathbb{T}^2)$ such that  $M = \phi N$. Further, perhaps we can replace $N$ with  $H^2(\mathbb{T}^2)$ to show that $M$ has Beurling-type.

\begin{lemma}
Let $M \subset H^2(\mathbb{T}^2)$ be shift-invariant. Suppose that for every $\varphi \in M$, we have $\varphi = \psi g$ for some inner function $\psi$. Let $N = \{g : \psi g \in M\}$ such that $M = \psi N$. Then set $N$ is shift-invariant.
\end{lemma}
\begin{proof}
Let  $N = \{g : \psi g \in M\}$. Then $zN = \{zg : \psi g \in M\}$. Since $\psi zg = z \psi g$ and $M$ is shift-invariant, we have $z \psi g \in M$ so that  $zg \in N$. That is $N$ is invariant under multiplication by $z$. Similarly, $N$ is invariant under multiplication by $w$ so that $N$ is shift-invariant. 
\end{proof}

\begin{theorem}
Let $M \subsetneq H^2(\mathbb{T}^2)$ be shift-invariant. $M$ is not full range if and only if $M = \psi N$ for some nonconstant inner function $\psi$ and some shift-invariant $N$.
\end{theorem}
\begin{proof}
    Assume $M = \psi N$ for some nonconstant inner function $\psi$ and some shift-invariant $N$. By Corollary 11 in [Mandrekar On Buerling type and equiv] $(M)_1 \subset (N)_1$   and $(M)_2 \subset (N)_2$. Also, 
    \begin{displaymath}
(M)_1 = \overline{span}
\{ \underset{n \geq 0}{\bigcup}{\overline{w}^n M} \} = \overline{span}
\{ \underset{n \geq 0}{\bigcup}{\overline{w}^n } \psi N \} = \psi  \overline{span}
\{ \underset{n \geq 0}{\bigcup}{\overline{w}^n N} \} = \psi (N)_1 
    \end{displaymath}
    \begin{displaymath}
(M)_2 = \overline{span}
\{ \underset{n \geq 0}{\bigcup}{\overline{z}^n M} \} = \overline{span}
\{ \underset{n \geq 0}{\bigcup}{\overline{z}^n } \psi N \} = \psi  \overline{span}
\{ \underset{n \geq 0}{\bigcup}{\overline{z}^n N} \} = \psi (N)_2 
    \end{displaymath}
Since $\psi$ is nonconstant inner, $\psi (z, w) =  \underset{i, j \geq 0}{\sum} c_{ij} z^i w^j$, where at least one of the coefficients in $\{c_{ij}: i \geq 1 $, or $j \geq 1\}$ is nonzero. Say $c_{i_0 j_0} \neq 0$ where $j_0 \geq 1$.  If $\psi (N)_2 = \mathcal{H}_2$, then $1 \in \psi (N)_2$ so that $\overline{\psi} \in (N)_2 \subset \mathcal{H}_2$. Since $\overline{\psi (z, w)} =  \underset{i, j \geq 0}{\sum} \overline{c_{ij}} \overline{z}^i \overline{w}^j$ where $c_{i_0 j_0} \neq 0$ and $j_0 \geq 1$, this implies that $\mathcal{H}_2$ contains an element with a negative power of $w$. Contradicting the definition of $\mathcal{H}_2$. 
\newline
\newline
unfinished**
If $\psi (N)_1 = \mathcal{H}_1$ then as $\psi (N)_1 \subset (N)_1$ %since $(N)_1$ is invariant under multiplication by $z, w, \text{and} \, \overline{w}$. 
this implies $ (M)_1 = (\psi N)_1 = \psi (N)_1 = (N)_1$ as $(N)_1 \subset \mathcal{H}_1$. Since $M = \psi N \subsetneq N$  However,  (show proper subset which implies $\mathcal{H}_1$ proper subset of $(N)_1$).
\end{proof}
\end{comment}
%However, in the $H^2(\mathbb{T}^2)$ setting every Beurling-type shift-invariant subspace has this property. This follows since by \cite{ACD86}, a shift-invariant subspace $M \subseteq H^2(\mathbb{T}^2)$ with finite codimension has full range (definition on pg 5 of \cite{ACD86}). By \cite{M88}, shift-invariant subspaces of the form $\phi H^2(\mathbb{T}^2)$, where $\phi(z, w)\in H^2(\mathbb{T}^2)$ is an inner function, do not have full range unless $\phi$ is constant. Therefore, a Beurling-type shift-invariant subspace, $\phi H^2(\mathbb{T}^2)$, cannot have finite codimension unless $\phi$ is constant. That is, shift-invariant subspaces of the form $\phi H^2(\mathbb{T}^2)$  have infinite codimension unless $\phi H^2(\mathbb{T}^2) = H^2(\mathbb{T}^2)$ (since $\phi H^2(\mathbb{T}^2) = H^2(\mathbb{T}^2)$ if and only if $\phi$ is constant). 

%In Theorem 3.6 in [ref: philipp] it was required that the inner function $\varphi(\xi)$ for which $Ker(V) = \varphi H^2(\mathbb{T})$ was not a finite Blaschke product. By Theorem 3.14 [ref: Invar Sub Radjavi], an invariant subspace $\varphi H^2(\mathbb{T})$ has finite codimension if and only if the inner function, $\varphi$, is a finite Blaschke product. Thus in Theorem 3.6, it was necessary that $\varphi H^2(\mathbb{T})$ had infinite codimension.  

%By [Argawal, Clark, Douglas] when  $\phi(z, w)$ is an inner function, any invariant subspace $M \subseteq H^2(\mathbb{T}^2)$ satisfies $ \phi M \subseteq M$ with equality if and only if  $\phi$ is constant. Also, due to %s stated in section 2 page 4 as well as in the proof of Corollary 3 
%By [ref: Argawal, Clark, Douglas], a shift-invariant subspace $M \subseteq H^2(\mathbb{T}^2)$ with finite codimension has full range (definition on pg 5 of ref: argawal douglas clark). By [ref: mandrekar], shift-invariant subspaces of the form $\phi H^2(\mathbb{T}^2)$, where $\phi(z, w)\in H^2(\mathbb{T}^2)$ is an inner function, do not have full range unless $\phi$ is constant. Thus, a Beurling-type shift-invariant subspace, $\phi H^2(\mathbb{T}^2)$, cannot have finite codimension unless $\phi$ is constant. That is, shift-invariant subspaces of the form $\phi H^2(\mathbb{T}^2)$  have infinite codimension unless $\phi H^2(\mathbb{T}^2) = H^2(\mathbb{T}^2)$ (since $\phi H^2(\mathbb{T}^2) = H^2(\mathbb{T}^2)$ if and only if $\phi$ is constant). 

%For example, it has been shown in [Nakazi '88] that every shift-invariant subspace, $M \subset H^2(\mathbb{T}^2)$, of finite co-dimension has the form $FN$ where $N$ is an invariant subspace of $L^2(\mathbb{T}^2)$ containing $H^2(\mathbb{T}^2)$ and $F$ is an inner function. 

\section{Frames by Unilateral Iterations of bounded Commuting Operators}
In what follows, we characterize the pairs of bounded commuting operators which can be used to generate frames for a separable infinite-dimensional Hilbert space. Critical for the characterization provided in Theorem \ref{char} is the following notion of similarity.
 \begin{definition} \label{similar}
Let $H, K$ be complex separable infinite-dimensional Hilbert Spaces. Let $T_1$ and $ T_2$ be commuting operators  in $B(H)$ and let  $V_1$ and $ V_2$ be commuting operators in $B(K)$. We say the triples $(T_1, T_2, \varphi)$ and $(V_1, V_2, f)$ are similar and write 
\newline
$(T_1, T_2, \varphi) \cong (V_1, V_2, f)$ if there exists an invertible map $L \in B(H, K)$ such that $LT_{1}L^{-1} = V_{1}$, $LT_{2}L^{-1} = V_{2}$, and $L\varphi = f$ where $\varphi \in H$ and $f \in K$.
\end{definition}
Note that the similarity relation in the above definition implies that $LT_{1}^iT_{2}^jL^{-1} = V_{1}^iV_{2}^j$ for all $i, j \geq 0$.


%Note that our congruence condition, $(T_1, T_2, \varphi)\cong (S_{\theta_1} S_{\theta_2}, P_{K_\theta} \mathbbm{1}_{\mathbb{T}^2})$ implies $LT_{1}L^{-1} = S_{\theta_1}$ and $LT_{2}L^{-1} = S_{\theta_2}$ for some $L \in GL(K_{\theta}, H)$. 

 %Given an invariant subspace $M \subset H^2(\mathbb{T}^2)$, let $E = M \ominus S_wM$. By Corollary 1 [Nakazi '88], $S_zE \subset E$ if and only if $M = FH^2(\mathbb{T}^2)$ where $F$ is inner. 

\begin{lemma} \label{similarity}
Suppose $(T_1, T_2, \varphi) \cong (V_1, V_2, f)$. Then $\{T_{1}^iT_{2}^j\varphi\}_{i,j \geq 0}$ is a frame if and only if $\{V_{1}^iV_{2}^jf\}_{i,j \geq 0}$ is a frame. In the affirmative case,  the operator $L$ in Definition \ref{similar} is unique. 
\end{lemma}
\begin{proof}
Since $\underset{i,j \geq 0}{\sum} |\langle V_{1}^iV_{2}^jf, g \rangle |^2 = \underset{i,j \geq 0}{\sum} |\langle LT_{1}^iT_{2}^jL^{-1} L\varphi , g\rangle |^2 = \underset{i,j \geq 0}{\sum} |\langle T_{1}^iT_{2}^j\varphi, L^*g \rangle |^2$, the statement follows since $L$ is a topological isomorphism so that $\{T_{1}^iT_{2}^j\varphi\}_{i,j \geq 0}$ being a frame implies $\{LT_{1}^iT_{2}^j\varphi\}_{i,j \geq 0}$ is a frame and thus $\{V_{1}^iV_{2}^jf\}_{i,j \geq 0}$ is a frame by the equation above. Also, assuming $\{V_{1}^iV_{2}^jf\}_{i,j \geq 0}$ is a frame implies $\{LT_{1}^iT_{2}^j\varphi\}_{i,j \geq 0}$ is a frame so that $\{T_{1}^iT_{2}^j\varphi\}_{i,j \geq 0} = \{L^{-1}LT_{1}^iT_{2}^j\varphi\}_{i,j \geq 0}$ is a frame as well. Uniqueness of $L$ holds as for any $f \in H$, we have $f = \underset{i,j \geq 0}{\sum} c_{i,j} T_1^iT_2^j\varphi$ so that 

\begin{displaymath}  
Lf = L \underset{i,j \geq 0}{\sum} c_{i,j} T_1^iT_2^j\varphi = \underset{i,j \geq 0}{\sum} c_{i,j} LT_1^iT_2^jL^{-1} L\varphi
= \underset{i,j \geq 0}{\sum} c_{i,j} V_1^iV_2^jf.
\end{displaymath}

\end{proof}

%Note that an invariant subspace of $H^2(\mathbb{T}^2)$ is invariant under multiplication by all elements of within  $\mathbb{C}[z, w]$. By [Yang Bidisc I] $H^2(\mathbb{T}^2) = clos(span\{z^iw^j : z, w \in \mathbb{T}, i,j \geq 0 \})$ 
%\cap clos\{\mathbb{C}[z, w]\} = \{f \in H^2(\mathbb{D}^2)$ : $f(z, w) = \underset{i,j \geq 0}{\sum} c_{ij} z^{i}w^{j}$, where $i, j \in \mathbb{N}_0\}. 


%Within the literature, determine useful conditions in $H^2(\mathbb{T}^2)$ setting to give Beurling type subspaces invariant under both $S_z$ and $S_w$. (Theorem 1 by Mandrekar seems to be very useful) Show equivalence similar to what is given in theorem 3.6 [Christensen Philipp]  between a frame $\{T_1^iT_2^j\varphi\}_{i,j \geq 0}$ for H where $T_1$ and $T_2$ (doubly commute? following theorem by Mandrekar) commute and a system generated by products of compressed $S_z$ and $S_w$ shifts in a ``model space''. Obtain similarity between $T_1$ and compressed $S_z$, and similarity between $T_2$ and compressed $S_w$ and obtain map $L(T_1^nT_2^m\varphi)=P_{\theta}z^nw^m \mathbbm{1}_{\mathbb{T}^2}$ where $P_{\theta}$ is the projection onto said model space. 
%Will $Ker(V)$ in $H^2(\mathbb{T}^2)$ setting satisfy condition in Mandrekar Theorem so that it is of Beurling type (or should we make it an assumption that it does)? 

 
 %\begin{definition}
 %Given a subspace $M \subset H$, the codimension of $M$ is $dim(H \ominus M)$, where $ 0 \leq dim(K) \leq \infty$ for any $K \subset H$.
 %\end{definition}
 The following lemma shows that the kernel of the synthesis operator for a frame generated by unilateral iterations of bounded commuting operators on a fixed vector must be an joint invariant subspace of the two right shift operators on $\ell^{2}(\mathbb{N}_0 \times \mathbb{N}_0)$.
\begin{lemma}
Suppose $\{T_1^iT_2^j\varphi\}_{i,j \geq 0}$ is a frame for H where $T_1, T_2  \in B(H)$ commute. Let $U: \ell^{2}(\mathbb{N}_0 \times \mathbb{N}_0) \to H$ be the synthesis operator of the frame. Let $R_1$ be the right shift in the first component, $R_2$ the right shift in the second component. Then $Ker(U)$ is invariant under $R_1$ and $R_2$.
\end{lemma}
\begin{proof}
Let $\{c_{i, j}\}_{i, j \geq 0} \in Ker(U)$. Then $U \{c_{i, j}\}_{i, j \geq 0} =  \underset{i,j \geq 0}{\sum} c_{i,j} T_1^iT_2^j\varphi = 0$.
\newline 
Note that
\begin{displaymath}
U (R_1 \{c_{i, j}\}_{i, j \geq 0}) =  \underset{i,j \geq 0}{\sum} c_{i,j} T_1^{i+1}T_2^j\varphi =  \underset{j \geq 0}{\sum} c_{0,j} T_1T_2^j\varphi + \underset{j \geq 0}{\sum} c_{1,j} T_1^{2}T_2^j\varphi + \ldots
\end{displaymath}
\begin{displaymath}
= T_1(  \underset{j \geq 0}{\sum} c_{0,j} T_2^j\varphi + \underset{j \geq 0}{\sum} c_{1,j} T_1T_2^j\varphi + \ldots) = T_1( \underset{i,j \geq 0}{\sum} c_{i,j} T_1^iT_2^j\varphi) = T_1(0) = 0.
\end{displaymath}
Thus $Ker(U)$ is invariant under $R_1$. Similarly, $Ker(U)$ is invariant under $R_2$
\end{proof}
%With this lemma we will have $Ker(U)$ corresponding to a shift-invariant subspace of $H^2(\mathbb{T}^2)$. 

%(Needed $Ker(U)$ invariant under unilateral shifts in Theorem 3.6 [Christensen Phillip] to obtain $Ker(V)$ invariant under multiplication by $z$.  By theorem 2.4 in [Christensen operator rep of sequences], for a frame $\{T^nf\}_{n \geq 0}$, $T$ bounded implies $Ker(U)$ is invariant under unilateral shifts. Here we need invariance under multiplication by $z$ and $w$ so needed to show boundedness of  both commuting operators gives invariance under both right shifts). 

%By [Rudin Polydiscs pg 121], given $\phi_1(z,w)$ and $\phi_2(z,w)$ are inner functions with the same zeros, $\phi_1(z,w)/\phi_2(z,w)$ is a constant inner function. 

%here just to keep track of papers 

%\begin{proposition}
%Let $\{T^n\varphi\}_{n \in \mathbb{N}_0}$ be an overcomplete frame for $H$. The inner function in Theorem 3.6 [Christensen Philipp] and the one given in Theorem 15 [Bailey] differ at most by a unimodular constant factor.
%\end{proposition}
%\begin{proof}
%By Proposition 2.3 [Christensen frame properties operators] given a frame has the form $\{T^n\varphi\}_{n \in \mathbb{N}_0}$, the frame is linearly independent. Thus due to the identification of  $H^2(\mathbb{T})$ with $H^2(\mathbb{D})$ we have equivalence of the statements in Theorem 3.6 [Christensen Phillip] and Theorem 15 [Bailey]. Let $W: \ell^2(\mathbb{N}_0) \to H^2(\mathbb{T})$ so that $W(c_0, c_1, \ldots) = \underset{n \geq 0}{\sum}c_n z^n$. Let $c \in Ker(U)$ be the scalar sequence given in Theorem 15 [Bailey] such that $Wc = \psi$ with $W(Ker(U)) = \psi  H^2(\mathbb{T})$. We aim to show that for $Ker(V) = Ker(U \mathcal{F}) = \phi H^2(\mathbb{T})$ we have $\phi H^2(\mathbb{T}) = \psi  H^2(\mathbb{T})$ so that the inner functions $\phi$ (given as $h$ in Theorem 3.6 [Chr Philipp]) and $\psi$ differ by a unimodular constant factor, where $\mathcal{F}: H^2(\mathbb{T}) \to \ell^2(\mathbb{N}_0)$ is the Fourier Transform. Let $f \in Ker(V)$. Since $f \in H^2(\mathbb{T})$ we have $f = \underset{n \geq 0}{\sum}a_n z^n$ for some $\{a_n\}_{n \in \mathbb{N}_0} \in \ell^2(\mathbb{N}_0)$. Note that $\mathcal{F} f = \{a_n\}_{n \in \mathbb{N}_0}$ so that since $0 = Vf = U \mathcal{F} f = U \{a_n\}_{n \in \mathbb{N}_0}$, we get $\{a_n\}_{n \in \mathbb{N}_0} \in Ker(U)$. Since $f = \underset{n \geq 0}{\sum}a_n z^n = W \{a_n\}_{n \in \mathbb{N}_0}$ we have $ f \in W(Ker(U))$. Now, let $f \in W(Ker(U))$. Then $f = \underset{n \geq 0}{\sum}c_n z^n $ for some $\{c_n\}_{n \in \mathbb{N}_0} \in Ker(U)$. Thus, $Vf = U \mathcal{F} f = U (\{c_n\}_{n \in \mathbb{N}_0}) = 0$. That is, $f \in Ker(V) $. Hence $\phi H^2(\mathbb{T}) = Ker(V) = W (Ker(U)) = \psi H^2(\mathbb{T})$. By Theorem 3.9 [Invar Radjavi] it follows that $\phi$ and $\psi$ differ at most by a unimodular constant factor.
%\end{proof}

%In Theorem 3.6 [Christensen Philipp] when the frame $\{T^n \varphi\}_{n \in \mathbb{N}_0}$ is a Riesz basis, the function $\phi$ ($h$ in their notation) is 0. Similarly in Theorem 15 [Bailey], since $Nul(U) = 0$ in this case, the sequence $c = 0$ so that $\psi = Wc = 0$. Thus the functions are equal in this case.

%By Yang [Bidisc I] we can identify functions in $H^2(\mathbb{D}^2)$ with those in $H^2(\mathbb{T}^2)$. 

Note that by \cite{R85}, the functions in $H^2(\mathbb{T}^2)$ are the radial limits of functions in $H^2(\mathbb{D}^2)$. By \cite{M88,R69} every $f \in H^2(\mathbb{D}^2)$ has the form $f(z_1, z_2) = \underset{i,j \geq 0}{\sum} c_{ij} z_{1}^{i}z_{2}^{j}$, where $i, j \in \mathbb{N}_0$ and $z_1, z_2 \in \mathbb{D}$  and the scalar sequence $\{c_{ij}\}_{i, j \geq 0}$ satisfies  $\underset{i,j \geq 0}{\sum}|c_{ij}|^2 < \infty$. Also, by \cite{M88} given $f \in H^2(\mathbb{D}^2)$ such that $f(z_1, z_2) = \underset{i,j \geq 0}{\sum} c_{ij} z_{1}^{i}z_{2}^{j}$, the radial limits of this function exist almost everywhere so that we can identify $f$ with $\tilde{f}(t_1, t_2) = \underset{i,j \geq 0}{\sum} c_{ij} t_{1}^{i}t_{2}^{j} \in H^2(\mathbb{T}^2)$, where $t_1, t_2 \in \mathbb{T}$. Thus all elements of $H^2(\mathbb{T}^2)$ can be written in the form $\underset{i,j \geq 0}{\sum} c_{ij} z^{i}w^{j}$ with $\{c_{ij}\}_{i, j \geq 0} \in \ell^2(\mathbb{N}_0 \times \mathbb{N}_0)$ and $z, w \in \mathbb{T}$.




Let $A: \ell^2(\mathbb{N}_0 \times \mathbb{N}_0) \to H^2(\mathbb{T}^2)$ where $A(\{c_{ij}\}_{i, j \geq 0}) = \underset{i,j \geq 0}{\sum} c_{ij} z^{i}w^{j}$. Then $A$ maps $\ell^2(\mathbb{N}_0 \times \mathbb{N}_0)$ onto $H^2(\mathbb{T}^2)$ and as for any $ f = \underset{i,j \geq 0}{\sum} a_{ij} z^{i}w^{j}  \in H^2(\mathbb{T}^2)$ we have 
\newline
$\|f\|_{H^2(\mathbb{T}^2)} = \|\{a_{ij}\}_{i, j \geq 0}\|_{\ell^2(\mathbb{N}_0 \times \mathbb{N}_0)}$, this map is unitary. 

Let $R_1$ and $R_2$ be the right shifts on $\ell^2(\mathbb{N}_0 \times \mathbb{N}_0)$ in the first and second components respectively. One can show that 
\newline
$AR_1 \{c_{ij}\}_{i, j \geq 0} = S_z A \{c_{ij}\}_{i, j \geq 0}$ and $AR_2 \{c_{ij}\}_{i, j \geq 0} = S_w A \{c_{ij}\}_{i, j \geq 0}$ so that $AR_1A^{-1} = S_z$ and $AR_2A^{-1} = S_w$. Also, as $A$ is unitary we have $AR_1^{*}A^{-1} = S_z^{*}$ and $AR_2^{*}A^{-1} = S_w^{*}$ as well.

%\begin{definition}
%A pair of isometries $V_1$ and $V_2$ doubly commute on a Hilbert Space, $H$, if $V_1$ commutes with $V_2$ and $V_1$ commutes with $V_2^{*}$.
%\end{definition}
%The following theorem by Mandrekar characterizes the Beurling-type shift-invariant subspaces of $H^2(\mathbb{T}^2)$.
%\begin{theorem}[Mandrekar]
%A non-trivial shift-invariant subspace $M \subset H^2(\mathbb{T}^2)$ is of the form $\phi H^2(\mathbb{T}^2)$ with $\phi(z,w)$ inner if and only if $S_z$ and $S_w$ doubly commute on $M$.
% \end{theorem}
 Applying the observations above, one can show that the images of joint invariant subspaces for $R_1$ and $R_2$ of $\ell^{2}(\mathbb{N}_0 \times \mathbb{N}_0)$ under the operator $A$ are exactly the submodules of $H^2(\mathbb{T}^2)$. Using the fact that the image under $A$ of the  kernel of the synthesis operator for a frame  of the form $\{T_{1}^iT_{2}^j\varphi\}_{i,j \geq 0}$ is then a submodule in $H^2(\mathbb{T}^2)$, assuming the right shift operators doubly commute on this space implies that the corresponding submodule is of Beurling-type.
\begin{lemma}\label{beur}
Let $\{T_{1}^iT_{2}^j\varphi\}_{i,j \geq 0}$ be an overcomplete frame for $H$ where $T_1, T_2 \in B(H)$ commute. Let $U: \ell^2(\mathbb{N}_0 \times \mathbb{N}_0) \to H$ be the synthesis operator for the frame. Assume that $R_1, R_2$ doubly commute on $Ker(U)$. Then the submodule $Ker(V) = Ker(U \mathcal{F}) \subset H^2(\mathbb{T}^2)$ is of the form $\phi H^2(\mathbb{T}^2)$, where $\phi(z,w)$ is an inner function. 
\end{lemma}
\begin{proof}
Suppose $R_1, R_2$ doubly commute on $Ker(U)$. As $S_z$ and $S_w$ commute on $H^2(\mathbb{T}^2)$, we want to show $S_z$ and $S_w^{*}$ commute on the invariant subspace $A(Ker(U))$. 
\newline
Given $R_1 R_2^{*} = R_2^{*} R_1$ on $Ker(U)$, we have on $A(Ker(U))$, 
\begin{displaymath}
S_zS_w^{*}= AR_1 A^{-1} A R_2^{*}A^{-1} =AR_1 R_2^{*}A^{-1} = AR_2^{*} R_1A^{-1} = AR_2^{*}A^{-1} A R_1A^{-1} =  S_w^{*} S_z.
\end{displaymath} Also, it holds that $A(Ker(U)) = Ker(U \mathcal{F})$ and therefore by \cite{M88}, $Ker(U \mathcal{F}) = \phi H^2(\mathbb{T}^2)$ for some inner function $\phi(z,w)$.
\end{proof}


% next step is to prove $\{z^nw^m\}_{m,n \geq 0} forms orthonormal basis for H^2(\mathbb{T}^2). Similar to proof in lemma 2.10 [Koca and Sahin BEURLING-TYPE BIDISC] show it is a complete orthonormal set. 
The following fact is well-known; however, we provide it here for clarity. 
\begin{lemma} \label{pars}

The set $\{S_{z}^n S_{w}^m {1}_{\mathbb{T}^2}\}_{m,n \geq 0} = \{z^nw^m\}_{m,n \geq 0}$  forms an orthonormal basis for $H^2(\mathbb{T}^2)$.
\end{lemma}
\begin{comment}
    
\begin{proof}
By \cite{Con90} we only need to show that $\{z^mw^n\}_{m,n \geq 0}$ is a complete orthonormal set in $H^2(\mathbb{T}^2)$. Since  $\{z^mw^n\}_{m,n \geq 0} \subset H^2(\mathbb{T}^2)$ we have 
\newline
$\overline{span}(\{z^mw^n\}_{m,n \geq 0}) \subseteq H^2(\mathbb{T}^2)$. By the above observations we have that  every element in $H^2(\mathbb{T}^2)$ can be written in the form $\underset{i,j \geq 0}{\sum} c_{j} z^{i}w^{j}$ with $\{c_{ij}\}_{i, j \geq 0} \in \ell^2(\mathbb{N}_0 \times \mathbb{N}_0)$. Hence it follows that $ H^2(\mathbb{T}^2) \subseteq \overline{span}(\{z^mw^n\}_{m,n \geq 0})$. That is, $ H^2(\mathbb{T}^2) = \overline{span}(\{z^mw^n\}_{m,n \geq 0})$ and $\{z^mw^n\}_{m,n \geq 0}$ is complete in $ H^2(\mathbb{T}^2)$. For any $m, n \geq 0$, 
\begin{displaymath}
\| z^m w^n \|^{2} = \langle z^m w^n , z^m w^n\rangle = \int_{\mathbb{T}^2}  z^m w^n \overline{z}^m \overline{w}^n d \mu 
\end{displaymath}
\begin{displaymath}
= \int_{\mathbb{T}^2}  z^m w^n {z}^{-m} {w}^{-n} d \mu = \int_{\mathbb{T}^2} 1 d \mu = \mu(\mathbb{T}^2) = 1. 
\end{displaymath}
Thus for any $m, n \geq 0$, $\| z^m w^n \|$ = 1. Let $(m, n) \neq (i, j)$ where $i, j \geq 0$ and $m, n \geq 0$. Then  
\begin{displaymath}
\langle z^m w^n , z^i w^j\rangle = \int_{\mathbb{T}^2}  z^m w^n \overline{z}^i \overline{w}^j d \mu = \int_{\mathbb{T}} (\int_{\mathbb{T}}  z^m w^n \overline{z}^i \overline{w}^j dz)dw 
\end{displaymath}
\begin{displaymath}
= (\int_{\mathbb{T}}  z^m \overline{z}^i dz) (\int_{\mathbb{T}} w^n \overline{w}^j dw)  = (\langle z^m , z^i \rangle_{H^2(\mathbb{T})})(\langle w^n , w^j\rangle_{H^2(\mathbb{T})}) = 0
\end{displaymath}
since $(m, n) \neq (i, j)$ and for $\xi \in \mathbb{T}$, $\{\xi^k\}_{k \geq 0}$ is orthonormal in $H^2(\mathbb{T})$ so that at least one of the terms in the product is 0.
\end{proof}
\end{comment}
\begin{lemma} \label{twovar}
Let $M \subset H^2(\mathbb{T}^2)$ be a submodule with infinite codimension. Set $K = H^2(\mathbb{T}^2) \ominus M$ and $S_{K_z} = P_{K} \left.S_{z}\right|_{K} $ and  $S_{K_w} = P_{K} \left.S_{w}\right|_{K}$. 
Then $S_{K_z}^m S_{K_w}^n P_{K} {1}_{\mathbb{T}^2} = P_{K} z^mw^n$ for all $m, n \geq 0$.
\end{lemma}
\begin{proof}
Let $f \in H^2(\mathbb{T}^2)$. Then  $P_{K} S_{z}^m S_{w}^n f = P_{K} S_{z}^m S_{w}^n P_{K} f + P_{K} S_{z}^m S_{w}^n P_{M} f = P_{K} S_{z}^m S_{w}^n P_{K} f$. So that for all $m, n \geq 0$ we have $P_{K} S_{z}^m S_{w}^n =P_{K} S_{z}^m S_{w}^n P_{K}$.
\newline
By the line above also, $P_{K} S_{z}^m  = P_{K} S_{z}^m P_{K} $ and $P_{K} S_{w}^n  = P_{K} S_{w}^n P_{K} $ from which it follows that on $K$, $P_{K} S_{z}^m P_{K} = S_{K_z}^m$ and $P_{K} S_{w}^n P_{K} = S_{K_w}^n$.

In particular, 
\begin{displaymath}
P_{K} z^m w^n =P_{K} S_{z}^m S_{w}^n {1}_{\mathbb{T}^2} = P_{K} S_{z}^m S_{w}^n P_{K} {1}_{\mathbb{T}^2} =
P_{K} S_{z}^m P_{K} (S_{w}^n P_{K} {1}_{\mathbb{T}^2})
\end{displaymath}
\begin{displaymath}
 = P_{K} S_{z}^m P_{K} P_{K} (S_{w}^n P_{K} {1}_{\mathbb{T}^2}) = S_{K_z}^m(P_{K} S_{w}^n P_{K} {1}_{\mathbb{T}^2}) =  S_{K_z}^m S_{K_w}^n P_{K} {1}_{\mathbb{T}^2}.
\end{displaymath}
\end{proof}

The following result shows that any frame given by forward iterations of a pair of bounded commuting operators on a single vector has the property that the operator pair must be similar to the two-variable Jordan block, that is the the pair  $(S_{K_z}, S_{K_w})$, on an infinite-dimensional quotient module $K$. 

\begin{theorem}\label{char}
Let $\{T_1^iT_2^j\varphi\}_{i,j \geq 0} \subset H$, where $T_1, T_2  \in B(H)$ commute. Then $\{T_1^iT_2^j\varphi\}_{i,j \geq 0}$ is an overcomplete frame if and only if there exists a nontrivial submodule $M \subset H^2(\mathbb{T}^2)$ with $dim(H^2(\mathbb{T}^2) \ominus M) = \infty$ such that $(T_1, T_2, \varphi)\cong (S_{K_z}, S_{K_w}, P_{K} {1}_{\mathbb{T}^2})$ where $P_{K}$ is the orthogonal projection onto the quotient module $K = H^2(\mathbb{T}^2) \ominus M$.  
\end{theorem}

\begin{proof}
Assume there exists a nontrivial invariant subspace $M \subset H^2(\mathbb{T}^2)$ with 
\newline
$dim(H^2(\mathbb{T}^2) \ominus M) = \infty$ such that $(T_1, T_2, \varphi)\cong (S_{K_z} S_{K_w}, P_{K}{1}_{\mathbb{T}^2})$ where $P_{K}$ is the orthogonal projection onto $K = H^2(\mathbb{T}^2) \ominus M$ and $S_{K_z} = P_{K} \left.S_{z}\right|_{K} $ and  $S_{K_w} = P_{K} \left.S_{w}\right|_{K}$. By Lemma \ref{pars} and Lemma \ref{twovar}, it follows that the system $\{S_{K_z}^m S_{K_w}^n P_{K} {1}_{\mathbb{T}^2}\}_{m, n \geq 0} = \{P_{K} z^mw^n\}_{m, n \geq 0}$  is an overcomplete Parseval frame for $K$. Hence, by %(uniqueness part of)
Lemma \ref{similarity}, we have  $\{T_1^iT_2^j\varphi\}_{i,j \geq 0}$ is an overcomplete frame for $H$. Conversely, suppose $\{T_1^iT_2^j\varphi\}_{i,j \geq 0}$ is an overcomplete frame for $H$. Define $V = U \mathcal{F}$. Then $V: H^2(\mathbb{T}^2) \to H$ is surjective so that $Ker(V)$ has infinite codimension. That is, $K = H^2(\mathbb{T}^2) \ominus Ker(V)$ satisfies $dim(K) =\infty$. Note that $V {1}_{\mathbb{T}^2} = \varphi$ since $\mathcal{F} {1}_{\mathbb{T}^2} $ gives the scalar sequence in $\ell^2(\mathbb{N}_0 \times \mathbb{N}_0) $ that has the value $1$ in the $(0, 0)$ component and the value $0$ everywhere else. Thus $ V {1}_{\mathbb{T}^2} = U \mathcal{F} {1}_{\mathbb{T}^2} = \varphi$. Similarly for each pair $(i, j)$, $VS_{z}^iS_{w}^j {1}_{\mathbb{T}^2} = Vz^iw^j = T_{1}^iT_{2}^j \varphi$. Set $W = \left.V\right|_{K}$. Then $\varphi = V {1}_{\mathbb{T}^2} = V(P_{K} {1}_{\mathbb{T}^2} + P_{Ker(V)} {1}_{\mathbb{T}^2}) = V P_{K} {1}_{\mathbb{T}^2} = W P_{K} {1}_{\mathbb{T}^2}$ and for every $i, j \geq 0$, $T_{1}^iT_{2}^j \varphi = Vz^iw^j = V P_{K} z^iw^j = W P_{K} z^iw^j= W S_{K_z}^i S_{K_w}^j P_{K} {1}_{\mathbb{T}^2}$. 
\newline
Note that $W \in B(K, H)$ is invertible and as for any $f \in H$, $f = \underset{i,j \geq 0}{\sum} c_{i,j} T_1^iT_2^j\varphi$ so that 
\begin{displaymath}
W S_{K_z}^m S_{K_w}^n W^{-1}f = W S_{K_z}^m S_{K_w}^n W^{-1} \underset{i,j \geq 0}{\sum} c_{i,j} T_1^iT_2^j\varphi = W S_{K_z}^m S_{K_w}^n  \underset{i,j \geq 0}{\sum} c_{i,j} W^{-1}T_1^iT_2^j\varphi 
\end{displaymath}
\begin{displaymath}
=W S_{K_z}^m S_{K_w}^n  \underset{i,j \geq 0}{\sum} c_{i,j} S_{K_z}^i S_{K_w}^j P_{K} {1}_{\mathbb{T}^2} = W \underset{i,j \geq 0}{\sum} c_{i,j} S_{K_z}^{i+m} S_{K_w}^{j+n} P_{K} {1}_{\mathbb{T}^2}  
\end{displaymath}
\begin{displaymath}
=\underset{i,j \geq 0}{\sum} c_{i,j} W S_{K_z}^{i+m} S_{K_w}^{j+n} P_{K} {1}_{\mathbb{T}^2} = \underset{i,j \geq 0}{\sum} c_{i,j} T_1^{i+m}T_2^{j+n}\varphi = T_1^{m}T_2^{n} \underset{i,j \geq 0}{\sum} c_{i,j} T_1^iT_2^j\varphi = T_1^{m}T_2^{n} f.
\end{displaymath}
\end{proof}

The following corollary provides a seamless extension from Theorem 3.6 in \cite{CHP20}. Applying Lemma \ref{beur} we impose a condition which ensures that the quotient module $K$ in Theorem \ref{char} is given by a submodule of Beurling type. This condition enables us to remove the assumption that the quotient module $K$ has infinite dimension in Theorem \ref{char} as by Proposition \ref{infcodim} proper submodules of Beurling-type always have infinite codimension. 
It follows directly from Lemma 1 in \cite{R85} that if for some inner functions $\phi_1$ and $\phi_2$ we have
$\phi_1 H^2(\mathbb{T}^2) = \phi_2 H^2(\mathbb{T}^2)$, then $\phi_1(z,w)/\phi_2(z,w)$ is a constant. 
%If  $\phi_1 H^2(\mathbb{T}^2) = \phi_2 H^2(\mathbb{T}^2)$, then the zero set of the functions $\phi_1$ and $\phi_2$ must differ at most by a set of measure $0$. Thus given $\phi_1 H^2(\mathbb{T}^2) = \phi_2 H^2(\mathbb{T}^2)$ we can assume $\phi_1$ and $\phi_2$ have the same zero set.
Hence when 
$\phi_1 H^2(\mathbb{T}^2) = \phi_2 H^2(\mathbb{T}^2)$, the inner functions $\phi_1$ and $\phi_2$ differ by a unimodular constant factor, $\phi_1(z,w)/\phi_2(z,w)$. From this, we have uniqueness of the inner function in the following corollary up to a unimodular constant factor.  
 \begin{corollary}\label{B-type}
Let $\{T_1^iT_2^j\varphi\}_{i,j \geq 0} \subset H$, where $T_1, T_2  \in B(H)$ commute, satisfy the property that the operators $R_1, R_2$ doubly commute on the kernel of the synthesis operator for $\{T_1^iT_2^j\varphi\}_{i,j \geq 0}$. Then $\{T_1^iT_2^j\varphi\}_{i,j \geq 0}$ is an overcomplete frame for $H$ if and only if there exists a proper Beurling-type submodule $\phi H^2(\mathbb{T}^2)$, with $\phi (z,w)$ a unique inner function, such that $(T_1, T_2, \varphi)\cong (S_{K_z}, S_{K_w}, P_{K} {1}_{\mathbb{T}^2})$ where $P_{K}$ is the orthogonal projection onto the quotient module $K = H^2(\mathbb{T}^2) \ominus \phi H^2(\mathbb{T}^2$).    
\end{corollary}
\begin{proof}
This theorem holds as a special case of Theorem \ref{char} where the shift-invariant subspace $M$ is of Beurling-type. By Proposition \ref{beur}, Beurling-type shift-invariant subspaces always have infinite codimension.

\end{proof}
%Proof following thm 57 here: Assume there exists a proper invariant subspace $\theta H^2(\mathbb{T}^2)$ with $\theta (z,w)$ inner such that $(T_1, T_2, \varphi)\cong (S_{\theta_1} S_{\theta_2}, P_{K_\theta} {1}_{\mathbb{T}^2})$ where $P_{K_\theta}$ is the orthogonal projection onto $K_\theta = H^2(\mathbb{T}^2) \ominus \theta H^2(\mathbb{T}^2)$ and $S_{\theta_1} = P_{K_\theta} \left.S_{z}\right|_{K_{\theta}} $ and  $S_{\theta_2} = P_{K_\theta} \left.S_{w}\right|_{K_{\theta}}$. By Lemma 11, it follows that the system $\{S_{\theta_1}^m S_{\theta_2}^n P_{K_\theta} {1}_{\mathbb{T}^2}\}_{m, n \geq 0} = \{P_{K_\theta} z^mw^n\}_{m, n \geq 0}$  is an overcomplete Parseval frame for $K_\theta$. Hence, by %(uniqueness part of)
%Lemma 2, we have  $\{T_1^iT_2^j\varphi\}_{i,j \geq 0}$ is an overcomplete frame for $H$. Conversely, suppose $\{T_1^iT_2^j\varphi\}_{i,j \geq 0}$ is an overcomplete frame for $H$. Define $V = U \mathcal{F}$. Then $V: H^2(\mathbb{T}^2) \to H$ and by Lemma 10, $Ker(V) = \theta H^2(\mathbb{T}^2)$ where $\theta(z, w)$ is an inner function. Note that $V {1}_{\mathbb{T}^2} = \varphi$ since $\mathcal{F} {1}_{\mathbb{T}^2} = $ gives the scalar sequence in $\ell^2(\mathbb{N}_0 \times \mathbb{N}_0) $ that has the value $1$ in the $(0, 0)$ component and the value $0$ everywhere else. Thus $ V {1}_{\mathbb{T}^2} = U \mathcal{F} {1}_{\mathbb{T}^2} = \varphi$. Similarly for each pair $(i, j)$, $VS_{z}^iS_{w}^j {1}_{\mathbb{T}^2} = Vz^iw^j = T_{1}^iT_{2}^j \varphi$. Set $W = \left.V\right|_{K_{\theta}}$. Then $\varphi = V {1}_{\mathbb{T}^2} = V(P_{K_\theta} {1}_{\mathbb{T}^2} + P_{Ker(V)} {1}_{\mathbb{T}^2}) = V P_{K_\theta} {1}_{\mathbb{T}^2} = W P_{K_\theta} {1}_{\mathbb{T}^2}$ and for every $i, j \geq 0$, $T_{1}^iT_{2}^j \varphi = Vz^iw^j = V P_{K_\theta} z^iw^j = W P_{K_\theta} z^iw^j= W S_{\theta_1}^i S_{\theta_2}^j P_{K_\theta} {1}_{\mathbb{T}^2}$. 
%\newline
%Note that $W \in GL(K_\theta, H)$ and as for any $f \in H$, $f = \underset{i,j \geq 0}{\sum} c_{i,j} T_1^iT_2^j\varphi$ so that 

%\begin{displaymath}
%W S_{\theta_1}^m S_{\theta_2}^n W^{-1}f = W S_{\theta_1}^m S_{\theta_2}^n W^{-1} \underset{i,j \geq 0}{\sum} c_{i,j} T_1^iT_2^j\varphi = W S_{\theta_1}^m S_{\theta_2}^n  \underset{i,j \geq 0}{\sum} c_{i,j} W^{-1}T_1^iT_2^j\varphi 
%\end{displaymath}
%\begin{displaymath}
%= W S_{\theta_1}^m S_{\theta_2}^n  \underset{i,j \geq 0}{\sum} c_{i,j} S_{\theta_1}^i S_{\theta_2}^j P_{K_\theta} {1}_{\mathbb{T}^2} = W \underset{i,j \geq 0}{\sum} c_{i,j} S_{\theta_1}^{i+m} S_{\theta_2}^{j+n} P_{K_\theta} {1}_{\mathbb{T}^2} \end{displaymath}
%\begin{displaymath}
%= \underset{i,j \geq 0}{\sum} c_{i,j} W S_{\theta_1}^{i+m} S_{\theta_2}^{j+n} P_{K_\theta} {1}_{\mathbb{T}^2} = \underset{i,j \geq 0}{\sum} c_{i,j} T_1^{i+m}T_2^{j+n}\varphi = T_1^{m}T_2^{n} \underset{i,j \geq 0}{\sum} c_{i,j} T_1^iT_2^j\varphi = T_1^{m}T_2^{n} f.
%\end{displaymath}

The following corollary provides a characterization of the pairs of bounded commuting operators which can be
used to generate a Riesz basis (or minimal frame) for a separable infinite-dimensional Hilbert space.
\begin{corollary}
Let $\{T_1^iT_2^j\varphi\}_{i,j \geq 0} \subset H$, where $T_1, T_2  \in B(H)$ commute. Then $\{T_1^iT_2^j\varphi\}_{i,j \geq 0}$ is a Riesz basis if and only if $(T_1, T_2, \varphi)\cong (S_z, S_w, {1}_{\mathbb{T}^2})$. 
\end{corollary}

\begin{proof}
Suppose  $(T_1, T_2, \varphi) \cong (S_z, S_w, {1}_{\mathbb{T}^2})$. We know that  $\{S_{z}^n S_{w}^m {1}_{\mathbb{T}^2}\}_{m,n \geq 0} = \{z^nw^m\}_{m,n \geq 0}$ \,  forms an orthonormal basis for $H^2(\mathbb{T}^2)$. It follows that there exists a topological isomorphism, $L$, such that $\{L S_{z}^n S_{w}^m {1}_{\mathbb{T}^2}\}_{n,m \geq 0} = \{T_1^nT_2^m\varphi\}_{n,m \geq 0}$. Since the image of an orthonormal basis under a topological isomorphism is a Riesz basis, we have  $\{T_1^iT_2^j\varphi\}_{i,j \geq 0}$ is a Riesz basis for $H$.
For the other direction, assume that $\{T_1^iT_2^j\varphi\}_{i,j \geq 0}$ is a Riesz basis. Define $V = U \mathcal{F}$. Then as $U$ is bijective, we get $V \in B(H^2(\mathbb{T}^2), H)$ is invertible and the result follows by setting $K = H^2(\mathbb{T}^2)$ in the proof of Theorem \ref{char} above.
\end{proof}
\section{Properties of Frames of Iterations via Commuting Operators}
The following proposition shows that for a frame of the form $\{T_1^iT_2^j\varphi\}_{i,j \geq 0}$ the adjoint of the iterating operator pair $(T^*_1, T^*_2)$ reflects the property as exhibited in the single operator case, that the adjoint of the  iterating operator and goes to zero in the strong operator topology \cite{AA17}. 


\begin{proposition}
    Let $\{T_1^iT_2^j\varphi\}_{i,j \geq 0} \subset H$ be a frame for $H$  where $T_1, T_2  \in B(H)$ commute. Then $(T_1^*)^i(T_2^*)^jf \to 0$ as $i, j \to \infty$.
\end{proposition}
\begin{proof}
    Note that 
    \begin{displaymath}
        (T_1^*)^i(T_2^*)^j = (T_2^jT_1^i)^* = (T_1^iT_2^j)^* = (T_2^*)^j(T_1^*)^i \, \, \, \forall i, j \geq 0.
        \end{displaymath}
Given $\{T_1^iT_2^j\varphi\}_{i,j \geq 0} \subset H$ is a frame with frame bounds $A , B >0, $ we have 
\begin{displaymath}
\underset{i, j \geq 0}{\sum}|\langle T_{1}^iT_{2}^j \, \varphi, (T_{1}^{m_1}T_{2}^{m_2})^* f \rangle |^2 =  \underset{i, j \geq 0}{\sum} |\langle T_{1}^{i+m_1}T_{2}^{j + m_2} \, \varphi, f \rangle |^2  =  \underset{j \geq m_2}{\sum} \, \, \underset{i \geq m_1}{\sum}|\langle T_{1}^{i}T_{2}^{j} \, \varphi, f \rangle |^2 \leq B \|f\|^2.
\end{displaymath}
Since the sum $\underset{j \geq m_2}{\sum} \, \, \underset{i \geq m_1}{\sum}|\langle T_{1}^{i}T_{2}^{j} \, \varphi, f \rangle |^2 $ converges, this implies that $|\langle T_{1}^{i}T_{2}^{j} \, \varphi, f \rangle |^2 \to 0$ as $m_1, m_2 \to \infty$. Also, as 
\begin{displaymath}
A \|(T_{1}^{m_1}T_{2}^{m_2})^* f\|^2 \leq \underset{i, j \geq 0}{\sum}|\langle T_{1}^iT_{2}^j \, \varphi, (T_{1}^{m_1}T_{2}^{m_2})^* f \rangle |^2  = \underset{j \geq m_2}{\sum} \, \, \underset{i \geq m_1}{\sum}|\langle T_{1}^{i}T_{2}^{j} \, \varphi, f \rangle |^2 
\end{displaymath}
we have $(T_1^*)^{m_1} (T_2^*)^{m_2}f = (T_{1}^{m_1}T_{2}^{m_2})^* f \to 0$ as as $m_1, m_2 \to \infty$
\end{proof}

In the case where the  frame $\{T_1^iT_2^j\varphi\}_{i,j \geq 0} \subset H$ satisfies the conditions of Lemma \ref{beur} we conjecture that the operator pair $(T_1, T_2)$ also reflects the property, as exhibited in the single operator case, that the iterating operator goes to zero in the strong operator topology \cite{CHP20}.

\begin{conjecture}
 Assume $\{T_1^iT_2^j\varphi\}_{i,j \geq 0}$ is an overcomplete frame for $H$, where $T_1, T_2  \in B(H)$ commute, such that that the operators $R_1, R_2$ doubly commute on the kernel of the synthesis operator  for the frame. Then for each $f \in H$ we have that  $T_1^iT_2^jf \to 0$ as $i, j \to \infty$.
\end{conjecture}
\begin{comment}
\begin{proposition}

\end{proposition}
\begin{proof}
Following Lemma C.2 in \cite{CHP20} it can be shown that for all $f \in H^2(\mathbb{T}^2)$, $P_{K_\theta} f = \theta \cdot P_{L^2(\mathbb{T}^2) \ominus H^2(\mathbb{T}^2)} (f \overline{\theta})$. Since for all $m, n \in \mathbb{N}_0$, $W S_{\theta_1}^m S_{\theta_2}^n W^{-1} = T_1^{m}T_2^{n}$ where $W$ is as defined in the proof of Theorem \ref{char}, it suffices to show that $S_{\theta_1}^m S_{\theta_2}^n g \to 0$ for all $g \in K_\theta$. Note that by the lemma above, we have $S_{\theta_1}^m S_{\theta_2}^n P_{K_\theta} {1}_{\mathbb{T}^2} = P_{K_\theta} z^m w^n {1}_{\mathbb{T}^2}$ which implies that for all $g \in K_\theta $, $S_{\theta_1}^m S_{\theta_2}^n g = P_{K_\theta} z^m w^n g$. Let $g \in K_\theta$, then 

$\|S_{\theta_1}^m S_{\theta_2}^n g\|_{L^2(\mathbb{T}^2)}^2 = \|P_{K_\theta} z^m w^n g\|_{L^2(\mathbb{T}^2)}^2 = \|\theta \cdot P_{L^2(\mathbb{T}^2) \ominus H^2(\mathbb{T}^2)} (z^m w^n g \overline{\theta})\|_{L^2(\mathbb{T}^2)}^2 $

$= \langle \theta \cdot P_{L^2(\mathbb{T}^2) \ominus H^2(\mathbb{T}^2)} (z^m w^n g \overline{\theta}), \theta \cdot P_{L^2(\mathbb{T}^2) \ominus H^2(\mathbb{T}^2)} (z^m w^n g \overline{\theta}) \rangle_{L^2(\mathbb{T}^2)} $

$= \langle \theta \overline{\theta} \cdot P_{L^2(\mathbb{T}^2) \ominus H^2(\mathbb{T}^2)} (z^m w^n g \overline{\theta}), P_{L^2(\mathbb{T}^2) \ominus H^2(\mathbb{T}^2)} (z^m w^n g \overline{\theta}) \rangle_{L^2(\mathbb{T}^2)}  $ 

$=\langle P_{L^2(\mathbb{T}^2) \ominus H^2(\mathbb{T}^2)} (z^m w^n g \overline{\theta}), P_{L^2(\mathbb{T}^2) \ominus H^2(\mathbb{T}^2)} (z^m w^n g \overline{\theta}) \rangle_{L^2(\mathbb{T}^2)}  $

$= \|P_{L^2(\mathbb{T}^2) \ominus H^2(\mathbb{T}^2)} (z^m w^n g \overline{\theta})\|_{L^2(\mathbb{T}^2)}^2 
$

$
= \underset{i, j >1} \sum | \langle z^m w^n g\overline{\theta}, z^{-i} w^{-j} \rangle_{L^2(\mathbb{T}^2)}|^2 $

$= \underset{i >m, j >n} \sum | \langle g \overline{\theta}, z^{-i} w^{-j} \rangle_{L^2(\mathbb{T}^2)}|^2$. Which tends to $0$ as $m, n \to \infty$.
\end{proof}
\end{comment}
Of interest in the area of Dynamical Frames (i.e. semigroup representation frames)  is determining semigroup representations (admitting frame vectors) for which their frame vectors are equivalent. Namely, when are each of the frame vectors for the frame representation of a semigroup the image of a fixed frame vector under an invertible operator in the commutant of the collection of operators given by the semigroup representation. Note that in the literature such frame representations are known as central frame representations \cite{BHKLL25}. 


Theorem 3.9 in \cite{CHP20} shows that the frame vectors are equivalent for any frame representation of the semigroup $Z_+$ (which are precisely the frames of the form $\{T^n f\}_{n\geq 0}$ where $T \in B(H)$). It was recently shown in \cite{BHKLL25} that the frame vectors are also equivalent for any frame representation of the semigroup $Z_+^n$. This result implies the following.

\begin{proposition}
   Let $\{T_1^iT_2^j \varphi\}_{i,j \geq 0} \subset H$, where $T_1, T_2  \in B(H)$ commute be a frame.  Then $\{T_1^iT_2^j f\}_{i,j \geq 0}$ is also a frame for $H$ if and only if $f = V \varphi$ for some invertible $V \in B(H)$ that commutes with both $T_1,$ and $T_2$.
   
\end{proposition}
\begin{proof}
    This follows immediately from Theorem 3.4 in \cite{BHKLL25}.
\end{proof}
\begin{proposition}
   Let $\{T_1^iT_2^j \varphi\}_{i,j \geq 0} \subset H$, where $T_1, T_2 \in B(H)$ commute be an overcomplete frame such that $R_1, R_2$ doubly commute on the kernel of its synthesis operator. Then for any frame $\{g_{i,j}\}_{i,j \geq 0} $  equivalent to $\{T_1^iT_2^j \varphi\}_{i,j \geq 0} \subset H$ (i.e. $ \{V(T_1^iT_2^j \varphi)\}_{i,j \geq 0} =  \{g_{i,j}\}_{i,j \geq 0}$ for some invertible $V \in B(H)$), we have $R_1, R_2$ doubly commute on the kernel of its synthesis operator as well. In particular, if $V$ commutes with both $T_1,$ and $T_2$, then replacing $\varphi$ with $g_{0 , 0}$ we obtain a frame of unilateral iterations satisfying the conditions of Corollary \ref{B-type}.
\end{proposition}

\begin{proof}
    Let $\{g_{i,j}\}_{i,j \geq 0} $ be a frame for $H$ (note that any frame can be indexed in this way). Observe that given the fact that  $\{g_{i,j}\}_{i,j \geq 0} $ is equivalent to $ \{T_1^iT_2^j \varphi\}_{i,j \geq 0}$,  the kernel of their synthesis operators coincide. This implies that $\{g_{i,j}\}_{i,j \geq 0} $ is an overcomplete frame and  $R_1, R_2$ doubly commute on the kernel of its  synthesis operator. The last statement follows as $g_{0 , 0}$ = $V\varphi$ so that $g_{0 , 0}$ is a frame vector for the pair $(T_1, T_2)$ and $\{g_{i,j}\}_{i,j \geq 0} = \{T_1^iT_2^j g_{0 , 0}\}_{i,j \geq 0}$.
\end{proof}

\begin{comment}
\section{Auxiliary Results}

 
\begin{lemma}
Suppose $H = H_1 \oplus H_2$ so that $H$ is the orthogonal direct sum of the Hilbert spaces $H_1$ and $H_2$.
Let $\{h_k\}_k \subset H_1$ be a frame for $H_1$ with frame bounds $A_1$ and $B_1$, and $\{g_l\}_l \subset H_2$ is a frame for $H_2$ with frame bounds $A_2$ and $B_2$. Then $\{h_k\}_k \cup \{g_l\}_l$ is a frame for $H$.
\end{lemma}
\begin{proof}
Let $f \in H$. Then $f = f_1 + f_2$ where $f_1 \in H_1$ and $f_2 \in H_2$ and
\begin{displaymath}
\sum_k |\langle f, h_k \rangle|^2 + \sum_l |\langle f, g_l \rangle|^2 = \sum_k |\langle f_1, h_k \rangle|^2 + \sum_l |\langle f_2, g_l \rangle|^2 \leq B_1 \|f_1\|^2 + B_2\|f_2\|^2 \leq B \|f\|^2
\end{displaymath}
as $f_1 \perp f_2$ where $B = max\{B_1, B_2\}$. Similarly the lower bound holds with lower frame bound $A = min\{A_1, A_2\}$.
\end{proof}

%By Rudin [function theory polydiscs], %pg 5 and pg 50
%if $f \in H^2(\mathbb{D}^2)$ is such that $f(z, w) = \underset{i,j \geq 0}{\sum} c_{ij} z^{i}w^{j}$, where $i, j \in \mathbb{N}_0$, then $\underset{i,j \geq 0}{\sum}|c_{ij}|^2 < \infty$. By Yang [Bidisc I] we can identify functions in $H^2(\mathbb{D}^2)$ with those in $H^2(\mathbb{T}^2)$. 

As in \cite{Y01}, let $H_w = H^2(\mathbb{T}^2) \ominus zH^2(\mathbb{T}^2) = \overline{span}\{w^j : j\geq 0\}\}$ and 
\newline
$H_z = H^2(\mathbb{T}^2) \ominus wH^2(\mathbb{T}^2) = \overline{span}\{z^j : j\geq 0\}\} $ where the closure is taken with respect to the $L^2$-norm.  Both $H_w$ and $H_z$ can be identified with $H^2(\mathbb{T})$.

\begin{proposition}
Let $H = H_1 \oplus H_2$ with $dim(H_1)=dim(H_2)= \infty$. Let $T_1, T_2 \in B(H)$. Then $\{T_{1}^nf_1\}_{n \in \mathbb{N}_0} \subset H_1$ is a frame for $H_1$ and $\{T_{2}^nf_2\}_{n \in \mathbb{N}_0} \subset H_2$ is a frame for $H_2$ so that  $\{T_{i}^nf_i\}_{n \in \mathbb{N}_0, i \in \{1, 2\}}$ is a frame for $H$
if and only if there exists a $S_z$ invariant subspace $M_1 \subset H_z$ and a $S_w$ invariant subspace $M_2 \subset H_w$  such that the pair $(T_1, f_1)$ is similar to $(S_{\theta}, P_{K_{\theta}}1_{\mathbb{T}^2})$ (where the compressed shift operator $S_{\theta} = P_{K_{\theta}} \left.S_z\right|_{K_{\theta}}$, with $K_\theta = H_z \ominus M_1$ and $P_{K_{\theta}}$ the orthogonal projection onto $K_{\theta}$) and the  pair $(T_2, f_2)$ is similar to $(S_{\gamma}, P_{K_{\gamma}}1_{\mathbb{T}^2})$ (where the compressed shift operator $S_{\gamma} = P_{K_{\gamma}} \left.S_w\right|_{K_{\gamma}}$, with $K_\gamma = H_w \ominus M_2$ and $P_{K_{\gamma}}$ the orthogonal projection onto $K_{\gamma}$). 
\end{proposition}

\begin{proof}
Suppose $\{T_{1}^nf_1\}_{n \in \mathbb{N}_0} \subset H_1$ is a frame for $H_1$ and $\{T_{2}^nf_2\}_{n \in \mathbb{N}_0} \subset H_2$ is a frame for $H_2$. Let  $\mathcal{F}_1 = \left.\mathcal{F}\right|_{H_z}$ where 
$\mathcal{F}: H^2(\mathbb{T}^2) \to  \ell^2(\mathbb{N}_0 \times \mathbb{N}_0)$ is the Fourier transform. For $f \in H_z$ each sequence $\{\hat{f}(j, l)\}_{j, l \geq 0} = (\mathcal{F}_1 f)(j, l)_{j, l \geq 0}$ satisfies $\hat{f}(j, l) = 0$ for all $l \neq 0$. Identifying the sequences $(\mathcal{F}_1 f)(j, l)_{j, l \geq 0} \in \ell^2(\mathbb{N}_0 \times \mathbb{N}_0)$ with sequences in $\ell^2(\mathbb{N}_0)$ by setting $(\tilde{\mathcal{F}_1} f)(j, l)_{j, l \geq 0}$ as the sequence $\{\tilde{f}(j)\}_{j \geq 0} = \{\hat{f}(j, 0)\}_{j \geq 0}$ set $V_1 = U_1 \tilde{\mathcal{F}_1}$ where $U_1$ is the synthesis operator for the sequence  $\{T_{1}^nf_1\}_{n \in \mathbb{N}_0}$. By Theorem 2.4 in \cite{CHS18}, since $T_1$ is bounded, $U_1$ is invariant under right shifts. It follows that $Ker(V_1)$ is invariant under $S_z$. Due to the identification between $H_z$ and $H^2(\mathbb{T})$, the conclusion of Theorem 3.6 in \cite{CHP20} holds by setting $H = H_1$, replacing $L^2_{+}$ with $H_z$, $h \in H_z$ an inner function that is not a finite Blaschke product.
\end{proof}
\end{comment}
%One of the most well known type of frames with various applications are generated by iterations of two bounded operators. These systems generated by iterations of translation and modulation operators on a single vector are called Gabor systems and numerous works over the span of many years have been focused on determining conditions for a given Gabor system to be a frame. In the case for Gabor systems, the translation and modulation operators do not commute. However, these operators \say{nearly} commute in the sense that interchanging the order of the translation and modulation operators can be achieved by multiplying by a modulus 1 constant factor. This fact motivates the following results. 
%\begin{lemma}
%Let $T_1 , T_2 \in B(H)$ and suppose $T_1 T_2 = q T_2 T_1$ where $q \in \mathbb{T}$ is a constant. Then $\{T_1^iT_2^j\varphi\}_{i,j \geq 0} = \{q^{ij}T_2^jT_1^i\varphi\}_{i,j \geq 0}$ is a frame if and only if $\{T_2^mT_1^n\varphi\}_{m,n \geq 0}$ is a frame.

%\end{lemma}

%\begin{proof}
%Let $T_1 , T_2 \in B(H)$ and suppose $T_1 T_2 = q T_2 T_1$ where $q \in \mathbb{T}$ is a constant. Assume $\{T_1^iT_2^j\varphi\}_{i,j \geq 0}$ is a frame. Since for any $f \in H$, $ \sum_{i, j \geq 0} |\langle f, T_1^iT_2^j\varphi \rangle|^2 = \sum_{i, j \geq 0} |\langle f, q^{ij}T_2^jT_1^i\varphi \rangle|^2 = \sum_{i, j \geq 0} |q^{ij}|^2 |\langle f, T_2^jT_1^i\varphi \rangle|^2 = \sum_{i, j \geq 0} |q|^{2ij} |\langle f, T_2^jT_1^i\varphi \rangle|^2 = \sum_{m, n \geq 0}  |\langle f, T_2^jT_1^i\varphi \rangle|^2 $
%\end{proof}

\begin{comment}
\section{Appendix}
 \begin{namedlemma}[Lemma C.2 in \cite{CHP20}]
Let $\phi$ be an inner function. Then
\begin{displaymath}
    P_{\phi(L^2(\mathbb{T}) \ominus H^2(\mathbb{T}))} H^2(\mathbb{T}) = \phi(L^2(\mathbb{T}) \ominus H^2(\mathbb{T})) \cap H^2(\mathbb{T}) = H^2(\mathbb{T}) \ominus \phi H^2(\mathbb{T}).
\end{displaymath}
Moreover,
\begin{displaymath}
     P_{\phi(L^2(\mathbb{T}) \ominus H^2(\mathbb{T})) \cap H^2(\mathbb{T})} f = \phi \cdot P_{(L^2(\mathbb{T}) \ominus H^2(\mathbb{T}))}(f \overline{\phi}), \, \forall f \in H^2(\mathbb{T}).
\end{displaymath}
\end{namedlemma}
 \begin{namedtheorem}[Theorem 2.4 in \cite{CHS18}]
Consider a sequence $\{f_k\}_{k=1}^{\infty}$ in $H$ which has a representation $\{f_k\}_{k=1}^{\infty} = \{T^n f_1\}_{n=0}^{\infty}$. for a linear operator  $T: \text{\text{span}} \{f_k\}_{k=1}^{\infty} \to \text{\text{span}} \{f_k\}_{k=1}^{\infty}$.  Then the following statements hold: 
\newline
$(i.)$ If $T$ is bounded, then the domain of the synthesis operator, $U$, and $Ker(U)$ are invariant under right shifts, and $\{\frac{ \|f_{k+1}\| }{\|f_k\|}\}_{k=1}^{\infty} \in \ell^{\infty}$ when $f_k \neq 0$ for all $k \in \mathbb{N}$ 
\newline
$(ii.)$  $T$ is bounded on $\text{\text{span}} \{f_k\}_{k=1}^{\infty}$ if and only if there is a positive constant $C$ so that 
    $\|UR \{c_k\}_{k=1}^{\infty}\| \leq C\|U \{c_k\}_{k=1}^{\infty}\|$
 for all finite sequences $\{c_k\}_{k=1}^{\infty}$.
\end{namedtheorem}
\end{comment}
\begin{thebibliography}{3}

\bibitem{ACD86} Agrawal, O.P., Clark, D., Douglas, R., Invariant subspaces in the polydisc, {\it Pacific Journal of Mathematics}, 121, 1-11 (1986).

\bibitem{ACCP22}
Aguilera A., Cabrelli, C., Carbajal, D., Paternostro, V.: Frames by orbits of two operators that commute. {\it Appl. Comp. Harm. Anal.} 66, 46-61 (2023).

\bibitem{ACMT17}                             
  Aldroubi, A., Cabrelli,  C., Molter, U.,  Tang, S.: Dynamical sampling. \textsl{Appl. Comput. Harmon. Anal.} 42, no. 3, 378–401 (2017)

 %\bibitem{ACM16} 
 %Aldroubi, A., Cabrelli,  C., Molter, U.,  Çakmak, A.F., Petrosyan, A., Iterative Actions of Normal Operators. {\it J. Funct. Anal.} 272, no. 3, 1121–1146, (2017).


 \bibitem{AA17}                         
 Aldroubi, A., Petrosyan, A., Dynamical sampling and systems from iterative actions of operators, Frames and other bases in abstract and function spaces {\it Appl. Numer. Harmon. Anal., Birkhäuser/Springer}, 15–26, (2017).

\bibitem{BHKLL25} Bailey, V., Han D., Kornelson, K., Larson, D., Liu, R., Dynamical frames and hyperinvariant subspaces. ArXiv:2505.19303, (2025).

\bibitem{B48} Beurling, A., On two problems concerning linear transformations in hilbert space {\it Acta Math.}, 81, 239–255 (1948).

\bibitem{CG03} Chen, X., Guo, K., Analytic hilbert modules {\it Chapman and Hall/CRC}, New York, (2003). 

%\bibitem{CH21} Christensen, O., Hasannasab, M., A survey on frame representations via dynamical sampling, Sampling, Approximation, and Signal Analysis, {\it Birkh\"auser} 349–370 (2024).

\bibitem{CHP20} Christensen, O., Hasannasab, M., Philipp, F., Frame properties of systems arising via iterative actions of operators {\it Appl. Comp. Harm. Anal.}, 46, 664-673 (2019).

%\bibitem{CHS18} Christensen, O. and Hasannasab, M. and Stoeva, D.T., Operator Representations of Sequences and Dynamical Sampling {\it Sampling Theory in Signal and Image Processing}, 17, (2018), 29-42.

\bibitem{Con90} Conway, J. B., A course in functional analysis {\it Springer-Verlag}, New York, (1990).

\bibitem{DP89} Douglas, R., Paulsen, V., Hilbert modules over function algebras, {\it Longman Scientific and Technical}, (1989).

\bibitem{Y00} Douglas, R., Yang, R., Operator theory in the hardy space over the bidisc, $I$ {\it Integr. equ. oper. theory.}, 38, 207–221 (2000).

\bibitem{Y01} Douglas, R., Yang, R., Operator theory in the hardy space over the bidisc, $III$ {\it J. Funct. Anal.}, 186, 521–545 (2001).

\bibitem{GMR16} Garcia, S.R., Mashreghi, J., Ross, W., Introduction to model spaces and their operators {\it Cambridge University Press}, Cambridge, (2016).

\bibitem{GM88} Ghatage, P., Mandrekar, V., On beurling type invariant subspaces of $L^2(\mathbb{T}^2)$ and their equivalence {\it Journal of Operator Theory}, 20, 83-89, (1988).

\bibitem{LYY11} Lu, Y., Yang, Y., Yang, R., An index formula for the two variable jordan block {\it Proceedings of the American Mathematical Society}, Vol. 139, No. 2, 511-520, (2011).

\bibitem{M88} Mandrekar, V., The validity of beurling theorems in polydiscs {\it Proceedings of the American Mathematical Society}, 103, 145-148 (1988).

\bibitem{MR07} Martínez-Avendaño, R., Rosenthal, P., An introduction to operators on the hardy-hilbert space.
{\it Graduate Texts in Mathematics, Springer}, New York, (2007).

\bibitem{RR03} Radjavi, H., Rosenthal, P., Invariant subspaces {\it Dover Publications, Inc.}, New York, (2003).
    
\bibitem{R69} Rudin, W., Function theory in polydiscs {\it W. A. Benjamin Inc.}, New York, (1969).

\bibitem{R85} Rudin, W., Invariant subspaces of $H^2$ on a torus, {\it J. Funct. Anal.}, 61, 378-384 (1985).

\bibitem{S15} Sarkar, J., An introduction to hilbert module approach to multivariable operator theory, Operator Theory, {\it Springer}, Basel,  (2015).

\bibitem{SD52} Schaeffer, A.C., Duffin, R.J., A class of nonharmonic fourier series.
\textsl{Trans. Amer. Math. Soc. 72, 341–366} (1952).

\bibitem{Y18} Yang, R., A brief survey on operator theory in $H^2(\mathbb{D}^2)$, Handbook of Analytic Operator Theory, {\it Chapman and Hall/CRC}, New York, (2019).

\end{thebibliography}
\end{document}

\bibitem{EHL94}C.A. Eschenbach, F.J. Hall, Z. Li,
          Eigenvalue frequency and consistent sign pattern matrices,
      {\it Czechoslovak Math. J.}, 44:3 (1994), 461--479.



\bibitem{EJ91} C.A. Eschenbach, C.R. Johnson,
          Sign patterns that require real, nonreal or pure imaginary eigenvalues,
      {\it Linear and Multilinear Algebra}, 29:3-4 (1991), 299--311.

\bibitem{EJ93} C.A. Eschenbach, C.R. Johnson,
          Sign patterns that require repeated eigenvalues,
      {\it Linear Algebra Appl.}, 190 (1993), 169--179.

\bibitem{Fang18} W. Fang, W. Gao, Y. Gao, F. Gong, G. Jing, Z. Li, Y. Shao,  L. Zhang, Minimum ranks of sign patterns and zero-nonzero patterns and point-hyperplane configurations, Linear Algebra Appl. 558 (2018), 44--62.

@article{VB22,
    title = {Frames via Operator Orbits: Structure of Subspaces Determining Boundedness},
    author = {Bailey, V.},
    journal = {arXiv:2201.09757},
    year = {2022}}

@article{Y18,
    title = {A brief survey on operator theory in $H^2(\mathbb{D}^2)$},
    author = {Yang, R.},
    journal = {arXiv:1810.00133},
    year = {2018}}


@article{Y01,
    title = {Operator Theory in the Hardy Space over the Bidisc, $III$},
    author = {Yang, R.},
    journal = {J. Funct. Anal.},
    volume =  {186},
    year = {2001},
    pages = {521–545}}


@book{BF01,
    title = {Modern Sampling Theory},
    author = {Benedetto, J. J. and Ferreira, P. J.S.G.},
    series = {Appl. Numer. Harmon. Anal.},
    publisher = {Birk\"auser Boston Inc.},
    location = {Boston, MA},
    year = {2001}}


@book{Chr03,
    author = {Christensen, O.},
    title = {An Introduction to Frames and Riesz Bases},
    series = {Appl. Numer. Harmon. Anal.},
    publisher = {Birk\"auser Boston Inc.},
    location = {Boston, MA},
    year = {2003}}


@article{CH19,
    title = {Frame properties of systems arising via iterative actions of operators},
    author = {Christensen, O. and Hasannasab, M.},
    journal = {Appl. Comp. Harm. Anal.},
    volume =  {46},
    year = {2019},
    pages = {664-673}}


@article{CH17,
    title = {Operator Representations of Frames: Boundedness, Duality, and Stability},
    author = {Christensen, O. and Hasannasab, M.},
    journal = {Integral Equations and Operator Theory},
    volume =  {88},
    year = {2017},
    pages = {483-499}}

@article{CHP20,
    title = {Frame Properties of Operator Orbits},
    author = {Christensen, O. and Hasannasab, M. and Philipp, F.},
    journal = {Math. Nach},
    volume =  {293},
    year = {2020},
    pages = {52-66}}


@article{CHS18,
    title = {Operator Representations of Sequences and Dynamical Sampling},
    author = {Christensen, O. and Hasannasab, M. and Stoeva, D.T.},
    journal = {Sampling Theory in Signal and Image Processing},
    volume =  {17},
    year = {2018},
    pages = {29-42}}

@article{CHR18,
    title = {Dynamical sampling and frame representations with bounded operators},
    author = {Christensen, O. and Hasannasab, M. and Rashidi, E.},
    journal = {SJ. Math. Anal. Appl.},
    volume =  {463},
    year = {2018},
    pages = {634–644}}

@article{CH21,
    title = {A survey on frame representations via dynamical sampling},
    author = {Christensen, O. and Hasannasab, M.},
    journal = {arXiv:2201.00038},
    year = {2021}}
    
@book{Con90,
    author = {Conway, J. B.},
    title = {A Course in Functional Analysis},
    edition = {Second Edition},
    publisher = {Springer-Verlag},
    location = {New York},
    year = {1990}}
    

@article{DS52,
    title = {A Class of Nonharmonic {F}ourier Series},
    author = {Duffin, R. J. and Schaeffer, A. C.},
    journal = {Trans. Amer. Math. Soc.},
    volume =  {72},
    year = {1952},
    pages = {341-366}}

        


@article{HL00,
    author = {Han, D. and Larson, D. R.},
    title = {Frames, Bases, and Group Representations},
    journal = {Mem. Amer. Math. Soc.},
    volume = {147},
    number =  {697},
    pages = {x+94},
    year = {2000}}
    
    

@book{Hei11,
    author = {Heil, C.},
    title = {A Basis Theory Primer},
    edition = {Expanded Edition},
    publisher = {Birkh\"auser},
    location = {Boston},
    year = {2011}}
    
    
@book{Hei18,
    author = {Heil, C.},
    title = {Metrics, Norms, Inner Products, and Operator Theory},
    series = {Appl. Numer. Harmon. Anal.},
    publisher = {Birkh\"auser/Springer},
    location = {Boston},
    year = {2018}}


@article{Hei07,
    title = {History and Evolution of the Density Theorem for Gabor Frames},
    author = {Heil, C.},
    journal = {J. Funct. Anal. Appl.},
    volume = {13},
    year = {2007},
    pages = {113–166}}  
    
@book{Hof88,
    author = {Hoffman, K.},
    title = {Banach Spaces of Analytic Functions},
    series = {Appl. Numer. Harmon. Anal.},
    publisher = {Dover Publications, Inc.},
    location = {New York},
    year = {1988}}

@book{HK71,
    author = {Hoffman, K. and Kunze, R.},
    title = {Linear Algebra},
    publisher = {Prentice-Hall, Inc.},
    location = {Englewoood Cliffs, N.J.},
    year = {1971}}
    
@article{IBL20,
    title = {Observing and tracking bandlimited graph processes from sampled measurements},
    author = {Isufi, E. and Banelli, P. and Lorenzo, P.D. and Leus, G.},
    journal = {Signal Processing},
    year = {2020},
    volume = {177}}

 @article{RBT19,
    title = {Fourier could be a data scientist: From graph Fourier transform to signal processing on graphs},
    author = {Ricaud, B. and Borgnat, P. and Tremblay N. and Gonçalves P. and  Vandergheynst, P.},
    journal = {Comptes Rendus Physique},
    year = {2019},
    volume = {20},   
    pages = {474-488}}
    
    
@article{LM09,
    title = {Spacial Super-resolution of a diffusion field by temporal oversampling in sensor networks},
    author = {Lu, Y. M. and Vetterli, M.},
    journal = {IEEE International Conference on Acoustics, Speech, and Signal Processing},
    year = {2009},
    pages = {2249-2252}}


@article{Pes08,
    author = {Pesenson, I.},
    title = {Sampling in Paley-Wiener Spaces on Combinatorial Graphs},
    journal = {Transactions of the American Mathematical Society},
    volume = {360},
    year = {2008},
    number = {10},
    pages = {5603-5627}}
    
    

@article{NS86,
    title = {Nonharmonic {F}ourier Series and its Applications},
    author = {Narukawa, K. and Suzuki, T.},
    journal = {Appl. Math. Optim.},
    volume = {14},
    year = {1986},
    number = {3},
    pages = {249-264}}
    


@book{SW99,
    author = {Schaefer, H. H. and Wolff, M. P.},
    title = {Topological Vector Spaces},
    edition = {Second Edition},
    series = {Graduate Texts in Mathematics},
    volume = {3},
    publisher = {Springer-Verlag},
    location = {New York},
    year = {1999}}
    
    

@book{SS78,
    author = {Steen, L. A. and Seebach Jr., J. A.},
    title = {Counterexamples in Topology},
    edition = {Second Edition},
    location = {New York-Heidelberg},
    publisher = {Springer-Verlag},
    year = {1978}}
    
    


@book{TJ89,
    title = {Banach-Mazur Distances and Finite-dimensional Operator Ideals},
    author = {Tomczak-Jaegermann, N.},
    series = {Pitman Monographs and Surveys in Pure and Applied Mathematics},
    volume = {38},
    location = {Harlow},
    publisher = {Longman Scientific and Technical},
    year = {1989}}
    
    

@book{Wal18,
    author = {Waldron, S.},
    title = {An Introduction to Finite Tight Frames},
    publisher = {Birkh\"auser},
    location = {Boston},
    year = {2018}}
    


@book{You01,
    title = {An Introduction to Nonharmonic Fourier Series},
    edition = {Revised First Edition},
    publisher = {Academic Press Inc.},
    location = {San Diego, CA},
    year = {2001},
    author = {Young, R. M.}}
    
    


@incollection{Zim01,
    title = {Normalized Tight Frames in Finite Dimensions},
    author = {Zimmermann, G.},
    pages = {249-252},
    booktitle = {Recent Progress in Multivariate Approximation (Witten-Bommerholz, 2000)},
    journal = {Internat. Ser. Numer. Math.},
    volume = {137},
    publisher = {Birkh\"auser},
    location = {Basel},
    year = {2001}}


@article{Breen22,
    title = {Minimum number of distinct eigenvalues allowed by a sign pattern},
    author = {Breen, J. and Brouwer, C. and Catral, M. and Cavers, M.  and van den Driessche, P. and VanderMeulen,  K.N.},
    journal = {Linear Algebra Appl.},
    volume = {654},
    year = {2022},
    pages = {311-338}}

@book{BS95,
    title = {Matrices of 
Sign-solvable Linear Systems},
    publisher = {Cambridge University Press},
    location = {Cambridge},
    year = {1995},
    author = {Brualdi, R.A. and Shader, B.L.}}

@article{Fang18,
    title = {Minimum ranks of sign patterns and zero-nonzero patterns and point-hyperplane configurations},
    author = {Fang, W. and Gao, W. and Gao, Y. and Gong, F. and Jing,  G. and Li, Z. and Shao, Y. and Zhang, L.},
    journal = {Linear Algebra Appl.},
    volume = {558},
    year = {2018},
    pages = {44-62}}

@article{DK89,
    title = {Low rank matrices with a given 
sign pattern},
    author = {Delsarte, P. and Kamp, Y.},
    journal = {SIAM J. Discrete Math.},
    volume = {2},
    year = {1989},
    pages = {51-63}}

@article{EHL94,
    title = {Eigenvalue frequency and consistent sign pattern matrices},
    author = {Eschenbach, C.A. and Hall, F.J and Li, Z.},
    journal = {Czechoslovak Math. J.},
    volume = {44:3},
    year = {1994},
    pages = {461-479}}

@article{EJ91,
    title = {Sign patterns that require real, nonreal or pure imaginary eigenvalues},
    author = {Eschenbach, C.A. and Johnson, C.R.},
    journal = {Linear and Multilinear Algebra},
    volume = {29:3-4},
    year = {1991},
    pages = {299-311}}

@article{EJ93,
    title = {Sign patterns that require repeated eigenvalues},
    author = {Eschenbach, C.A. and Johnson, C.R.},
    journal = {Linear Algebra Appl.},
    volume = {190},
    year = {1993},
    pages = {169-179}}

@article{Feng20,
    title = {Rank conditions for sign patterns that allow diagonalizability},
    author = {Feng, X.-L. and Gao, W. and Hall, F.J. and Jing,  G. and Li, Z. and Zagrodny, C. and Zhou, J.},
    journal = {Discrete Mathematics},
    volume = {343.5},
    year = {2020}}

 @incollection{HL14,
    title = {Sign pattern matrices},
    author = {Hall, F.J and Li, Z.},
    chapter = {42},
    booktitle = {Handbook of Linear Algebra},
    journal = {L. Hogben, ed.},
    publisher = {CRC Press},
    location = {Boca Raton, FL},
    year = {2014}}

@book{Horn13,
    title = {Matrix Analysis, 2nd ed.},
    publisher = {Cambridge University Press},
    location = {Cambridge},
    year = {2013},
    author = {Horn, R.A. and Johnson, C.R.}}

@article{HL99,
    title = {The moment and Gram matrices, distinct eigenvalues and zeroes, and rational criteria for diagonalizability},
    author = {Horn, R.A. and Lopatin, A.K.},
    journal = {Linear Algebra Appl.},
    volume = {299},
    year = {1999},
    pages = {153--163}}


@book{Sam47,
    title = {Foundations of Economic Analysis},
    publisher = {Harvard University Press},
    location = {Cambridge},
    year = {1947},
    author = {Samuelson, P.A.}}

@article{KP97,
    title = {Nonlocal sensitivity analysis of the eigensystem of a matrix with distinct eigenvalues},
    author = {Konstantinov, M.M. and Petkov, P.H. and Gu, D.W. and Postlethwaite, I. },
    journal = {Numerical Functional Analysis and Optimization},
    volume = {3-4},
    year = {1997},
    pages = {367-382}}

@article{LH02,
    title = {Sign patterns that require all distinct eigenvalues},
    author = {Li, Z. and Harris, L. },
    journal = {JP J. Algebra Number Theory Appl.},
    volume = {2:2},
    year = {2002},
    pages = {161-179}}

@article{LW88,
    title = {First- and second-order eigensensitivity of matrices with distinct eigenvalues},
    author = {Ling, Y.L and Wang, B.C. },
    journal = {International Journal of Systems Science},
    volume = {7},
    year = {1988},
    pages = {1053-1067}}

@article{Mar08,
    title = {Positive Polynomials and Sums of Squares},
    author = {Marshal, M.},
    journal = {American Math. Society},
    year = {2008}}


@book{Pras10,
    title = {Polynomials},
    publisher = {Springer-Verlag},
    location = {Berlin},
    year = {2010},
    author = {Prasolov, V.V.}}
    
@article{SG03,
    title = {Sign patterns that allow diagonalizability},
    author = { Shao, Y.L. and  Gao, Y. B. },
    journal = {Linear Algebra Appl.},
    volume = {359 (1-3)},
    year = {2003},
    pages = {113-119}}

@article{C96,
    title = {Frames Containing a Riesz Basis and Approximation of the Frame Coefficients Using Finite-Dimensional Methods},
    author = {Christensen, O.},
    journal = {J. Math. Anal. Appl.},
    volume = {199},
    year = {1996},
    pages = {256–270}}    

@article{CA02,
    title = {Decomposition of Riesz Frames and Wavelets into a Finite Union of Linearly Independent Sets},
    author = {Christensen, O. and Lindner, A.},
    journal = {Linear Algebra Appl.},
    volume = {355},
    year = {2002},
    pages = {147–159}}        

@article{ACMT17,
    title = {Dynamical Sampling},
    author = {Aldroubi, A. and Cabrelli,  C. and Molter, U. and  Tang, S.},
    journal = {Appl. Comput. Harmon. Anal.},
    volume = {42},
    year = {2017},
    pages = {378–401}}  

@book{MR07,
    title = {An Introduction to Operators on the Hardy-Hilbert Space},
    publisher = {Springer},
    location = {New York},
    year = {2007},
    author = {Martínez-Avendaño, R. and Rosenthal, P.}}

@article{B48,
    title = {On Two Problems Concerning Linear Transformations in Hilbert Space},
    author = {Beurling, A.},
    journal = {Acta Math.},
    volume = {81},
    year = {1948},
    pages = {239–255}}  

@book{RR03,
    title = {Invariant Subspaces},
    publisher = {Dover Publications, Inc.},
    location = {New York},
    year = {2003},
    author = {Radjavi, H. and Rosenthal, P.}}

 @incollection{AA17,
    title = {Dynamical Sampling and Systems from Iterative Actions of Operators},
    author = {Aldroubi, A. and Petrosyan, A.},
    booktitle = {Frames and other bases in abstract and function spaces},
    journal = {Appl. Numer. Harmon. Anal.},
    publisher = {Birkhäuser/Springer},
    year = {2017}}

@article{AHK18,
    title = {Dynamical Sampling with Additive Random Noise},
    author = {Aldroubi, A. and Huang, L. and Krishtal, I. and Ledeczi, A. and  Lederman, R.R. and Volgyesi, P. },
    journal = {Sampling Theory in Signal and Image Processing},
    volume = {17},
    year = {2018},
    pages = {153–182}}  


@article{ACMÇ17,
    title = {Iterative Actions of Normal Operators},
    author = {Aldroubi, A. and Cabrelli,  C. and Molter, U. and Çakmak, A.F. and Petrosyan, A.},
    journal = {J. Funct. Anal.},
    volume = {272},
    year = {2017},
    pages = {1121–1146}}  

@article{M88,
    title = {The Validity of Beurling Theorems in Polydiscs},
    author = {Mandrekar, V.},
    journal = {Proceedings of the American Mathematical Society},
    volume = {103},
    year = {1988},
    pages = {145-148}}  

@article{ACD86,
    title = {Invariant subspaces in the polydisc},
    author = {Agrawal, O.P. and Clark, D. and Douglas, R.},
    journal = {Pacific Journal of Mathematics},
    volume = {121},
    year = {1986},
    pages = {1-11}}  

@article{R85,
    title = {Invariant subspaces of $H^2$ on a torus},
    author = {Rudin, W.},
    journal = {J. Funct. Anal.},
    volume = {61},
    year = {1985},
    pages = {378-384}} 
    
@book{R87,
    title = {Real and Complex Analysis},
    publisher = {McGraw-Hill Book Co.},
    location = {New York},
    year = {1987},
    author = {Rudin, W.}}

@book{R91,
    title = {Functional Analysis},
    publisher = {McGraw-Hill, Inc.},
    location = {New York},
    year = {1991},
    author = {Rudin, W.}}


@book{R69,
    title = {Functional Theory in Polydiscs},
    publisher = {W. A. Benjamin Inc.},
    location = {New York},
    year = {1969},
    author = {Rudin, W.}}

@book{GMR16,
    title = {Introduction to Model Spaces and their Operators},
    publisher = {Cambridge University Press},
    location = {Cambridge},
    year = {2016},
    author = {Garcia, S.R. and Mashreghi, J. and Ross, W.}}

%\bibitem [1]{AB02}   
 %Abakumov, E., Borichev, A.: Shift Invariant Subspaces with Arbitrary Indices in $l^p$ Spaces. \textsl{J. Funct. Anal. 188, no. 1, 1–26} (2002)
%\bibitem [15] {HD58}                                   
%Helson, H., Lowdenslager, D.: Prediction Theory and Fourier Series in Several Variables. \textsl{Acta Math. 99, 165–202} (1958)

%\bibitem [17] {L59}                                   
%Lax, P.: Translation Invariant Spaces, \textsl{Acta Math. 101, 163–178} (1959)

